% Options for packages loaded elsewhere
\PassOptionsToPackage{unicode}{hyperref}
\PassOptionsToPackage{hyphens}{url}
\documentclass[
]{article}
\usepackage{xcolor}
\usepackage{amsmath,amssymb}
\setcounter{secnumdepth}{-\maxdimen} % remove section numbering
\usepackage{iftex}
\ifPDFTeX
  \usepackage[T1]{fontenc}
  \usepackage[utf8]{inputenc}
  \usepackage{textcomp} % provide euro and other symbols
\else % if luatex or xetex
  \usepackage{unicode-math} % this also loads fontspec
  \defaultfontfeatures{Scale=MatchLowercase}
  \defaultfontfeatures[\rmfamily]{Ligatures=TeX,Scale=1}
\fi
\usepackage{lmodern}
\ifPDFTeX\else
  % xetex/luatex font selection
\fi
% Use upquote if available, for straight quotes in verbatim environments
\IfFileExists{upquote.sty}{\usepackage{upquote}}{}
\IfFileExists{microtype.sty}{% use microtype if available
  \usepackage[]{microtype}
  \UseMicrotypeSet[protrusion]{basicmath} % disable protrusion for tt fonts
}{}
\makeatletter
\@ifundefined{KOMAClassName}{% if non-KOMA class
  \IfFileExists{parskip.sty}{%
    \usepackage{parskip}
  }{% else
    \setlength{\parindent}{0pt}
    \setlength{\parskip}{6pt plus 2pt minus 1pt}}
}{% if KOMA class
  \KOMAoptions{parskip=half}}
\makeatother
\usepackage{color}
\usepackage{fancyvrb}
\newcommand{\VerbBar}{|}
\newcommand{\VERB}{\Verb[commandchars=\\\{\}]}
\DefineVerbatimEnvironment{Highlighting}{Verbatim}{commandchars=\\\{\}}
% Add ',fontsize=\small' for more characters per line
\newenvironment{Shaded}{}{}
\newcommand{\AlertTok}[1]{\textcolor[rgb]{1.00,0.00,0.00}{\textbf{#1}}}
\newcommand{\AnnotationTok}[1]{\textcolor[rgb]{0.38,0.63,0.69}{\textbf{\textit{#1}}}}
\newcommand{\AttributeTok}[1]{\textcolor[rgb]{0.49,0.56,0.16}{#1}}
\newcommand{\BaseNTok}[1]{\textcolor[rgb]{0.25,0.63,0.44}{#1}}
\newcommand{\BuiltInTok}[1]{\textcolor[rgb]{0.00,0.50,0.00}{#1}}
\newcommand{\CharTok}[1]{\textcolor[rgb]{0.25,0.44,0.63}{#1}}
\newcommand{\CommentTok}[1]{\textcolor[rgb]{0.38,0.63,0.69}{\textit{#1}}}
\newcommand{\CommentVarTok}[1]{\textcolor[rgb]{0.38,0.63,0.69}{\textbf{\textit{#1}}}}
\newcommand{\ConstantTok}[1]{\textcolor[rgb]{0.53,0.00,0.00}{#1}}
\newcommand{\ControlFlowTok}[1]{\textcolor[rgb]{0.00,0.44,0.13}{\textbf{#1}}}
\newcommand{\DataTypeTok}[1]{\textcolor[rgb]{0.56,0.13,0.00}{#1}}
\newcommand{\DecValTok}[1]{\textcolor[rgb]{0.25,0.63,0.44}{#1}}
\newcommand{\DocumentationTok}[1]{\textcolor[rgb]{0.73,0.13,0.13}{\textit{#1}}}
\newcommand{\ErrorTok}[1]{\textcolor[rgb]{1.00,0.00,0.00}{\textbf{#1}}}
\newcommand{\ExtensionTok}[1]{#1}
\newcommand{\FloatTok}[1]{\textcolor[rgb]{0.25,0.63,0.44}{#1}}
\newcommand{\FunctionTok}[1]{\textcolor[rgb]{0.02,0.16,0.49}{#1}}
\newcommand{\ImportTok}[1]{\textcolor[rgb]{0.00,0.50,0.00}{\textbf{#1}}}
\newcommand{\InformationTok}[1]{\textcolor[rgb]{0.38,0.63,0.69}{\textbf{\textit{#1}}}}
\newcommand{\KeywordTok}[1]{\textcolor[rgb]{0.00,0.44,0.13}{\textbf{#1}}}
\newcommand{\NormalTok}[1]{#1}
\newcommand{\OperatorTok}[1]{\textcolor[rgb]{0.40,0.40,0.40}{#1}}
\newcommand{\OtherTok}[1]{\textcolor[rgb]{0.00,0.44,0.13}{#1}}
\newcommand{\PreprocessorTok}[1]{\textcolor[rgb]{0.74,0.48,0.00}{#1}}
\newcommand{\RegionMarkerTok}[1]{#1}
\newcommand{\SpecialCharTok}[1]{\textcolor[rgb]{0.25,0.44,0.63}{#1}}
\newcommand{\SpecialStringTok}[1]{\textcolor[rgb]{0.73,0.40,0.53}{#1}}
\newcommand{\StringTok}[1]{\textcolor[rgb]{0.25,0.44,0.63}{#1}}
\newcommand{\VariableTok}[1]{\textcolor[rgb]{0.10,0.09,0.49}{#1}}
\newcommand{\VerbatimStringTok}[1]{\textcolor[rgb]{0.25,0.44,0.63}{#1}}
\newcommand{\WarningTok}[1]{\textcolor[rgb]{0.38,0.63,0.69}{\textbf{\textit{#1}}}}
\usepackage{longtable,booktabs,array}
\newcounter{none} % for unnumbered tables
\usepackage{calc} % for calculating minipage widths
% Correct order of tables after \paragraph or \subparagraph
\usepackage{etoolbox}
\makeatletter
\patchcmd\longtable{\par}{\if@noskipsec\mbox{}\fi\par}{}{}
\makeatother
% Allow footnotes in longtable head/foot
\IfFileExists{footnotehyper.sty}{\usepackage{footnotehyper}}{\usepackage{footnote}}
\makesavenoteenv{longtable}
\setlength{\emergencystretch}{3em} % prevent overfull lines
\providecommand{\tightlist}{%
  \setlength{\itemsep}{0pt}\setlength{\parskip}{0pt}}
\usepackage{bookmark}
\IfFileExists{xurl.sty}{\usepackage{xurl}}{} % add URL line breaks if available
\urlstyle{same}
\hypersetup{
  hidelinks,
  pdfcreator={LaTeX via pandoc}}

\author{}
\date{}

\begin{document}

\section{OpenClaw Memory System: Architecture \& Technical Deep
Dive}\label{openclaw-memory-system-architecture-technical-deep-dive}

\textbf{nox-mem v3.7 --- March 2026, §5 empirical evaluation added May
2026 (Wave A G5 V3)}

\textbf{Author:} Luiz Antonio Busnello (Toto) \textbf{Platform:}
OpenClaw Autonomous Agent Platform \textbf{Infrastructure:} Hostinger
KVM4, Tailscale VPN, Debian Linux

\begin{center}\rule{0.5\linewidth}{0.5pt}\end{center}

\subsection{Abstract}\label{abstract}

This paper presents the architecture, implementation, and operational
characteristics of nox-mem, a persistent memory system designed for
autonomous AI agent fleets. The system provides hybrid search (combining
BM25 full-text, semantic vector similarity, and Reciprocal Rank Fusion),
an LLM-powered knowledge graph with temporal decay, cross-agent
intelligence sharing, and automated consolidation pipelines. Deployed in
production since March 14, 2026, the system manages 1,481 memory chunks
across 7 databases, 384 knowledge graph entities with 529 relations, and
serves 6 specialized AI agents with isolated yet interconnectable memory
spaces.

\begin{center}\rule{0.5\linewidth}{0.5pt}\end{center}

\subsection{1. Introduction}\label{introduction}

\subsubsection{1.1 Problem Statement}\label{problem-statement}

Large Language Model (LLM) agents operating in production environments
face a fundamental limitation: context window ephemerality. When a
conversation ends or context is compacted, the agent loses accumulated
knowledge, decisions, and operational state. For multi-agent systems
where specialized agents collaborate on complex tasks, this problem
compounds --- agents cannot learn from each other's experiences, cannot
recall past decisions, and cannot build institutional knowledge over
time.

\subsubsection{1.2 Design Goals}\label{design-goals}

nox-mem was designed with four core objectives:

\begin{enumerate}
\def\labelenumi{\arabic{enumi}.}
\tightlist
\item
  \textbf{Persistent Memory}: Survive context window resets, session
  boundaries, and agent restarts
\item
  \textbf{Intelligent Retrieval}: Return semantically relevant results,
  not just keyword matches
\item
  \textbf{Cross-Agent Intelligence}: Enable knowledge sharing across
  isolated agent workspaces
\item
  \textbf{Operational Autonomy}: Self-maintain through automated
  consolidation, pruning, and indexing
\end{enumerate}

\subsubsection{1.3 Scope}\label{scope}

The system operates within the OpenClaw platform, serving 6 AI agents
(Nox, Atlas, Boris, Cipher, Forge, Lex) on a single VPS with 4 vCPUs and
8GB RAM. Each agent has a distinct role and memory profile. The
workspace (shared memory) and individual agent databases form a
federated memory architecture.

\begin{center}\rule{0.5\linewidth}{0.5pt}\end{center}

\subsection{2. System Architecture}\label{system-architecture}

\subsubsection{2.1 Infrastructure
Overview}\label{infrastructure-overview}

The system runs on a Hostinger KVM4 VPS accessible via Tailscale VPN
(IP: 100.87.8.44). Five systemd-managed services provide the runtime
environment:

{\def\LTcaptype{none} % do not increment counter
\begin{longtable}[]{@{}llll@{}}
\toprule\noalign{}
Service & Port & Type & Function \\
\midrule\noalign{}
\endhead
\bottomrule\noalign{}
\endlastfoot
openclaw-gateway & 18789 & WebSocket & Agent communication gateway \\
nox-mem-watcher & --- & inotifywait & Filesystem event monitor \\
nox-mem-api & 18800 & HTTP/JSON & Dashboard data API \\
ollama & 11434 & HTTP & Local LLM inference (llama3.2:3b) \\
tailscaled & --- & WireGuard & VPN mesh connectivity \\
\end{longtable}
}

\subsubsection{2.2 Database Schema}\label{database-schema}

The primary storage is SQLite 3 with WAL (Write-Ahead Logging) mode for
concurrent access. The schema (version 3) contains:

\textbf{Core Tables:}

\begin{itemize}
\tightlist
\item
  \texttt{chunks} --- Memory fragments with full-text indexing

  \begin{itemize}
  \tightlist
  \item
    \texttt{id} (INTEGER PK), \texttt{source\_file} (TEXT),
    \texttt{chunk\_text} (TEXT), \texttt{chunk\_type} (TEXT)
  \item
    \texttt{source\_date} (TEXT), \texttt{is\_consolidated} (INTEGER),
    \texttt{memory\_type} (TEXT)
  \item
    \texttt{created\_at}, \texttt{updated\_at} (TEXT, ISO 8601)
  \item
    \texttt{metadata} (TEXT, JSON)
  \end{itemize}
\item
  \texttt{chunks\_fts} --- FTS5 virtual table with porter unicode61
  tokenizer

  \begin{itemize}
  \tightlist
  \item
    Content-sync triggers (INSERT, UPDATE, DELETE) maintain index
    consistency
  \item
    BM25 ranking with configurable column weights (1.0, 0.5, 0.5)
  \end{itemize}
\item
  \texttt{consolidated\_files} --- Processing state tracker

  \begin{itemize}
  \tightlist
  \item
    \texttt{source\_file} (TEXT PK), \texttt{status} (INTEGER:
    0=pending, 1=done, -1=failed)
  \end{itemize}
\item
  \texttt{meta} --- Key-value configuration store (schema\_version,
  cursors, metrics)
\end{itemize}

\textbf{Knowledge Graph Tables:}

\begin{itemize}
\tightlist
\item
  \texttt{kg\_entities} --- Named entities with type classification and
  mention counting

  \begin{itemize}
  \tightlist
  \item
    UNIQUE constraint on (name, entity\_type)
  \item
    TTL tracking via \texttt{first\_seen}, \texttt{last\_seen}
  \end{itemize}
\item
  \texttt{kg\_relations} --- Typed relationships between entities

  \begin{itemize}
  \tightlist
  \item
    Confidence scoring (0.0-1.0) with temporal decay
  \item
    TTL via \texttt{expires\_at} (90-day default),
    \texttt{last\_confirmed}
  \item
    Evidence linking via \texttt{evidence\_chunk\_id}
  \end{itemize}
\item
  \texttt{decision\_versions} --- Architectural decision version history

  \begin{itemize}
  \tightlist
  \item
    Supersession chain via \texttt{is\_current} flag and
    \texttt{superseded\_at} timestamp
  \end{itemize}
\end{itemize}

\textbf{Vector Tables (sqlite-vec):}

\begin{itemize}
\item
  \texttt{vec\_chunks} --- Virtual table storing float32 embeddings
  (3072 dimensions)
\item
  \texttt{vec\_chunk\_map} --- Rowid-to-chunk\_id mapping (sqlite-vec
  requires rowid-based access)
\item
  \texttt{dedup\_log} --- Suppressed duplicate tracking for audit
\end{itemize}

\subsubsection{2.3 Chunk Type Taxonomy}\label{chunk-type-taxonomy}

Memory chunks are classified into 10 types based on source file path
patterns:

{\def\LTcaptype{none} % do not increment counter
\begin{longtable}[]{@{}
  >{\raggedright\arraybackslash}p{(\linewidth - 6\tabcolsep) * \real{0.1333}}
  >{\raggedright\arraybackslash}p{(\linewidth - 6\tabcolsep) * \real{0.3333}}
  >{\raggedright\arraybackslash}p{(\linewidth - 6\tabcolsep) * \real{0.3333}}
  >{\raggedright\arraybackslash}p{(\linewidth - 6\tabcolsep) * \real{0.2000}}@{}}
\toprule\noalign{}
\begin{minipage}[b]{\linewidth}\raggedright
Type
\end{minipage} & \begin{minipage}[b]{\linewidth}\raggedright
Source Pattern
\end{minipage} & \begin{minipage}[b]{\linewidth}\raggedright
Current Count
\end{minipage} & \begin{minipage}[b]{\linewidth}\raggedright
Purpose
\end{minipage} \\
\midrule\noalign{}
\endhead
\bottomrule\noalign{}
\endlastfoot
team & \texttt{shared/} & 499 & Shared team knowledge, cross-agent
docs \\
daily & \texttt{memory/YYYY-MM-DD} & 161 & Daily operational notes \\
other & (default) & 126 & Unclassified content \\
decision & \texttt{memory/decisions.md} & 34 & Architectural and
strategic decisions \\
lesson & \texttt{memory/lessons.md} & 21 & Errors, corrections,
learnings \\
project & \texttt{memory/projects.md} & 11 & Active project tracking \\
pending & \texttt{memory/pending.md} & 8 & Incomplete tasks and
blockers \\
feedback & \texttt{memory/feedback/} & 6 & User and system feedback \\
person & \texttt{memory/people.md} & 6 & People profiles and contacts \\
digest & \texttt{memory/digests/} & 2 & Weekly summary reports \\
\end{longtable}
}

\subsubsection{2.4 Multi-Agent Memory
Architecture}\label{multi-agent-memory-architecture}

Each of the 6 agents operates with an isolated database at
\texttt{/root/.openclaw/agents/\{name\}/tools/nox-mem/nox-mem.db}. The
\texttt{OPENCLAW\_WORKSPACE} environment variable controls path
resolution across all modules, enabling the same nox-mem binary to
operate on different databases depending on the calling context.

\textbf{Agent Memory Distribution (as of March 23, 2026):}

{\def\LTcaptype{none} % do not increment counter
\begin{longtable}[]{@{}lllll@{}}
\toprule\noalign{}
Agent & Role & Chunks & DB Size & Dominant Type \\
\midrule\noalign{}
\endhead
\bottomrule\noalign{}
\endlastfoot
Nox & Chief of Staff & 185 & 268 KB & daily (91) \\
Boris & Head of Communications & 148 & 268 KB & team (50) \\
Forge & Code Reviewer & 182 & 292 KB & daily (135) \\
Atlas & Research & 30 & 128 KB & other (17) \\
Cipher & Security & 31 & 132 KB & other (12) \\
Lex & Legal/Compliance & 31 & 132 KB & other (12) \\
\textbf{Workspace} & \textbf{Shared} & \textbf{874} & \textbf{25.2 MB} &
\textbf{team (499)} \\
\end{longtable}
}

Total system memory: 1,481 chunks across 7 databases.

\begin{center}\rule{0.5\linewidth}{0.5pt}\end{center}

\subsection{3. Memory Pipeline}\label{memory-pipeline}

\subsubsection{3.1 Ingestion}\label{ingestion}

Files created or modified in monitored directories trigger the
inotifywait-based watcher service. The watcher implements:

\begin{itemize}
\tightlist
\item
  \textbf{Debounce logic}: 2-second delay to batch rapid successive
  writes
\item
  \textbf{File filtering}: Only \texttt{.md} and \texttt{.json} files
  are processed
\item
  \textbf{Recursion prevention}: \texttt{MEMORY.md} and
  \texttt{SESSION-STATE.md} are excluded to avoid feedback loops
\item
  \textbf{Heartbeat}: Touch \texttt{/tmp/nox-mem-watcher-heartbeat} on
  every event for liveness monitoring
\end{itemize}

Upon trigger, \texttt{ingestFile()} executes:

\begin{enumerate}
\def\labelenumi{\arabic{enumi}.}
\tightlist
\item
  Read file content with UTF-8 sanitization (fixes common mojibake
  patterns for Portuguese text)
\item
  Detect chunk type from relative file path
\item
  Extract date from filename pattern (YYYY-MM-DD)
\item
  Split content into semantic chunks:

  \begin{itemize}
  \tightlist
  \item
    Markdown: Split on H2/H3 headers, with sub-splitting for chunks
    exceeding 500 words
  \item
    JSON: Array items become individual chunks; object entries become
    key-value pairs
  \item
    Small chunks (\textless20 words) are merged with the previous chunk
  \end{itemize}
\item
  Delete existing chunks for the same source file (idempotent
  re-ingestion)
\item
  Insert new chunks via prepared statement transaction
\item
  Auto-vectorize if GEMINI\_API\_KEY is available (up to 20 chunks per
  file)
\end{enumerate}

\subsubsection{3.2 Consolidation}\label{consolidation}

Nightly consolidation (23:00, 5-minute stagger across agents) processes
daily notes into structured topic files:

\begin{enumerate}
\def\labelenumi{\arabic{enumi}.}
\tightlist
\item
  \textbf{Reindex}: Scan all \texttt{.md}/\texttt{.json} files in memory
  directories, rebuild chunk index
\item
  \textbf{Extract}: Use Ollama llama3.2:3b to identify facts, decisions,
  lessons, and action items from daily notes
\item
  \textbf{Append}: Add extracted content to topic files (decisions.md,
  lessons.md, people.md, projects.md, pending.md)
\item
  \textbf{Notion Sync}: Push structured items to ``Memoria \& Decisoes''
  Notion database (best-effort, non-blocking)
\item
  \textbf{Git Commit}: Auto-commit memory changes with standardized
  message format
\item
  \textbf{Session Update}: Refresh SESSION-STATE.md with current
  statistics
\end{enumerate}

\subsubsection{3.3 Deduplication}\label{deduplication}

Before insertion, chunks are checked for duplicates using a two-tier
strategy:

\begin{itemize}
\tightlist
\item
  \textbf{Primary}: Gemini cosine similarity with 0.85 threshold (when
  embeddings are available)
\item
  \textbf{Fallback}: Keyword overlap calculation with 60\% threshold
\item
  \textbf{Audit}: Suppressed duplicates are logged to
  \texttt{dedup\_log} table with reason and preview
\end{itemize}

\begin{center}\rule{0.5\linewidth}{0.5pt}\end{center}

\subsection{4. Hybrid Search System}\label{hybrid-search-system}

\subsubsection{4.1 Architecture}\label{architecture}

Search combines three complementary retrieval methods:

\textbf{Layer 1 --- FTS5 BM25 (Keyword)}

SQLite FTS5 with porter unicode61 tokenizer provides fast keyword
matching. Results are scored using BM25 with column weights
(chunk\_text: 1.0, source\_file: 0.5, chunk\_type: 0.5). Post-retrieval
boosting applies:

\begin{itemize}
\tightlist
\item
  Type boost: \texttt{decision} and \texttt{lesson} chunks receive 2.0x
  multiplier (higher signal-to-noise ratio)
\item
  Recency boost: Chunks from the last 7 days receive 1.5x multiplier
\end{itemize}

The query sanitizer strips special characters but preserves hyphens for
compound terms (e.g., ``nox-mem'').

\textbf{Layer 2 --- Gemini Semantic (Vector)}

Each chunk is embedded using Google's gemini-embedding-001 model (3072
dimensions) with task type RETRIEVAL\_DOCUMENT. Query embeddings use
task type RETRIEVAL\_QUERY for asymmetric similarity optimization.

Vectors are stored in sqlite-vec virtual tables. Retrieval uses cosine
distance with a map table (vec\_chunk\_map) bridging vec\_chunks rowids
to chunks.id values due to sqlite-vec's rowid-only constraint.

Scoring normalizes distances to a 0-10 scale with type and recency
boosting (1.5x and 1.2x respectively, lower than FTS5 to avoid
double-boosting in fusion).

\textbf{Layer 3 --- Reciprocal Rank Fusion (RRF)}

FTS5 and semantic results are merged using RRF with k=60:

\begin{verbatim}
RRF_score(d) = Σ 1/(k + rank_i(d))
\end{verbatim}

Documents appearing in both result sets receive combined scores, marked
as \texttt{match\_type:\ "hybrid"}. Content-prefix deduplication (first
50 characters) prevents near-duplicate results.

\subsubsection{4.2 Performance
Characteristics}\label{performance-characteristics}

The hybrid approach provides significant quality improvements over
single-method search:

{\def\LTcaptype{none} % do not increment counter
\begin{longtable}[]{@{}
  >{\raggedright\arraybackslash}p{(\linewidth - 6\tabcolsep) * \real{0.1944}}
  >{\raggedright\arraybackslash}p{(\linewidth - 6\tabcolsep) * \real{0.3056}}
  >{\raggedright\arraybackslash}p{(\linewidth - 6\tabcolsep) * \real{0.2222}}
  >{\raggedright\arraybackslash}p{(\linewidth - 6\tabcolsep) * \real{0.2778}}@{}}
\toprule\noalign{}
\begin{minipage}[b]{\linewidth}\raggedright
Query
\end{minipage} & \begin{minipage}[b]{\linewidth}\raggedright
FTS5 Only
\end{minipage} & \begin{minipage}[b]{\linewidth}\raggedright
Hybrid
\end{minipage} & \begin{minipage}[b]{\linewidth}\raggedright
Analysis
\end{minipage} \\
\midrule\noalign{}
\endhead
\bottomrule\noalign{}
\endlastfoot
``qual o proximo passo'' & 0 results & ROADMAP + PHASE-3 & Semantic
captures intent without keyword match \\
``nox-mem'' & 0 results & decisions.md + docs & Vector bypasses
tokenizer hyphen issues \\
``quem e o Toto'' & people.md & people.md + TEAM\_MEMORY & RRF combines
exact match + semantic context \\
\end{longtable}
}

\subsubsection{4.3 Cross-Agent Search}\label{cross-agent-search}

The \texttt{crossSearch()} function opens all 7 databases in read-only
mode, executes FTS5 queries in each, and merges results with agent
attribution. Deduplication uses content-prefix comparison to handle
shared documents that appear across multiple agent databases.

\begin{center}\rule{0.5\linewidth}{0.5pt}\end{center}

\subsection{5. Empirical Evaluation --- Wave A G5 V3 (May
2026)}\label{empirical-evaluation-wave-a-g5-v3-may-2026}

\begin{quote}
\textbf{Headline (canonical, 2026-05-19):} A8 full stack with active
salience reaches \textbf{nDCG@10 = 0.6237} on the entity-flavored golden
set (n=100), a \textbf{+78.8\% relative improvement over the G3 baseline
(0.3488)} measured prior to Wave A deployment, and \textbf{+9.4\% over
the mid-deployment G4 checkpoint (0.5702)}. The peak ablation isolating
\texttt{section\_boost} alone (A3) reaches 0.6228, recovering
\textbf{99.85\% of the full stack} --- section-aware ranking is the
dominant driver of the lift.
\end{quote}

\subsubsection{5.1 Setup}\label{setup}

The evaluation uses an entity-flavored golden set of 100 queries
(\texttt{entity-eval.db}), curated from production usage to exercise the
V10 schema's \texttt{section} and \texttt{pain} dimensions.
Configurations are toggled via environment-variable feature gates
(\texttt{NOX\_SALIENCE\_MODE}, \texttt{NOX\_DISABLE\_TIER\_BOOST},
\texttt{NOX\_ENABLE\_TIER\_BOOST},
\texttt{NOX\_DISABLE\_SECTION\_BOOST}, etc.), allowing isolation of
individual ranking components without code changes between runs. All
measurements occur post-deployment of PRs \#150 (salience formula +
tier\_boost off-by-default), \#151 (source\_type backfill of 67,949
chunks), and \#153 (search wiring). Reported nDCG@10 follows the
standard TREC formulation (gain by relevance, log-position discount).

Progression vs prior ablation generations:

{\def\LTcaptype{none} % do not increment counter
\begin{longtable}[]{@{}
  >{\raggedright\arraybackslash}p{(\linewidth - 8\tabcolsep) * \real{0.2000}}
  >{\raggedright\arraybackslash}p{(\linewidth - 8\tabcolsep) * \real{0.2000}}
  >{\raggedright\arraybackslash}p{(\linewidth - 8\tabcolsep) * \real{0.2000}}
  >{\raggedright\arraybackslash}p{(\linewidth - 8\tabcolsep) * \real{0.2000}}
  >{\raggedright\arraybackslash}p{(\linewidth - 8\tabcolsep) * \real{0.2000}}@{}}
\toprule\noalign{}
\begin{minipage}[b]{\linewidth}\raggedright
Generation
\end{minipage} & \begin{minipage}[b]{\linewidth}\raggedright
Date
\end{minipage} & \begin{minipage}[b]{\linewidth}\raggedright
A8 nDCG@10
\end{minipage} & \begin{minipage}[b]{\linewidth}\raggedright
Δ vs G3 baseline
\end{minipage} & \begin{minipage}[b]{\linewidth}\raggedright
Notes
\end{minipage} \\
\midrule\noalign{}
\endhead
\bottomrule\noalign{}
\endlastfoot
G3 baseline (pre-Wave A) & 2026-05-15 & 0.3488 & --- & Multiplicative
salience, tier\_boost on, section\_boost only via legacy code path \\
G4 mid-deployment & 2026-05-18 & 0.5702 & +63.5\% & Salience aditivo
wired but \texttt{active\ \textless{}\ shadow} puzzle observed \\
\textbf{G5 V3 canonical} & \textbf{2026-05-19} & \textbf{0.6237} &
\textbf{+78.8\%} & Wave A fully deployed; reversal
\texttt{active\ \textgreater{}\ shadow} cravado \\
\end{longtable}
}

The four sub-claims below decompose the +78.8\% total into measurable
contributions; the full G5 V3 matrix (12 configurations) is archived in
\texttt{audits/} and HANDOFF.md (\texttt{\#g5-v3-matrix-2026-05-19}).

\subsubsection{5.2 Claim 1 --- Additive salience outperforms
multiplicative}\label{claim-1-additive-salience-outperforms-multiplicative}

The Wave A formula replaces the legacy multiplicative
\texttt{salience\ =\ recency\ ×\ pain\ ×\ importance} with a
weighted-additive form (PR \#150):

\begin{verbatim}
salience = W_IMPORTANCE·importance + W_RECENCY·recency + W_PAIN·pain + W_ACCESS·access_score
W_IMPORTANCE = 0.55   W_RECENCY = 0.15   W_PAIN = 0.10   W_ACCESS = 0.20
\end{verbatim}

\textbf{Result.} With \texttt{NOX\_SALIENCE\_MODE=active}, A8 reaches
0.6237 vs.~0.6155 with \texttt{shadow} (A7) --- a +1.3\% lift and the
reversal of the G4 puzzle, where shadow had outranked active. The
multiplicative form concentrated 99.7\% of chunks in the {[}0.05,
0.40{]} salience range, dominated by 90.67\% of chunks at the default
\texttt{pain\ =\ 0.2} and 99.76\% of chunks with
\texttt{recency\ ∈\ {[}7,\ 30{]}} days; small differences in any factor
were swallowed by the product. The additive form exposes each dimension
proportionally to its calibrated weight, preserving signal from pain
spikes and importance heuristics without requiring all three factors to
be simultaneously non-default.

\subsubsection{\texorpdfstring{5.3 Claim 2 --- \texttt{section\_boost}
is the moat (99.85\% of the
gain)}{5.3 Claim 2 --- section\_boost is the moat (99.85\% of the gain)}}\label{claim-2-section_boost-is-the-moat-99.85-of-the-gain}

Isolating \texttt{section\_boost} alone (A3 ablation: section enabled,
tier off, source\_type off, salience shadow) yields \textbf{nDCG@10 =
0.6228 = 99.85\% of A8's full-stack 0.6237}. The V10 schema multipliers
--- \texttt{compiled\ =\ 2.0}, \texttt{frontmatter\ =\ 1.5},
\texttt{timeline\ =\ 0.8}, legacy = 1.0 --- together with the
entity-file format introduced in v3.7 (769 entity files × 3 sections ≈
2,307 boost-bearing chunks) explain the majority of the headline
improvement.

The negative control A11 (full stack minus \texttt{section\_boost})
drops to 0.5646, \textbf{−9.5\% relative to A8}, confirming the
contribution is not redundant with semantic embeddings or RRF fusion.
This is the architectural pivot the paper's narrative rests on:
section-aware boosting over an entity-file canonical form is the
load-bearing component, not the multiplicative salience formula that the
v1 paper draft over-emphasized.

\subsubsection{\texorpdfstring{5.4 Claim 3 --- \texttt{tier\_boost}
off-by-default is the correct
calibration}{5.4 Claim 3 --- tier\_boost off-by-default is the correct calibration}}\label{claim-3-tier_boost-off-by-default-is-the-correct-calibration}

Isolated, \texttt{tier\_boost} (boost for \texttt{chunks} flagged as
\texttt{tier=\textquotesingle{}core\textquotesingle{}}) is actively
harmful: A6 (tier only, no other boosts) reaches 0.4059, \textbf{−21\%
versus the no-boost baseline 0.5126}. Even integrated into the full
stack, A9 (full + tier enabled) drops to 0.5884, \textbf{−5.7\% versus
A8}. Inspection of the corpus reveals the cause:
\texttt{tier=\textquotesingle{}core\textquotesingle{}} chunks account
for only 3.96\% of the corpus and consist of memory-system internals
(lifecycle docs, schema metadata, operational runbooks) rather than user
content --- over-promoting them displaces directly-relevant entity
facts.

PR \#150 therefore makes tier\_boost \textbf{off by default} via
\texttt{NOX\_DISABLE\_TIER\_BOOST=1}, with an explicit
\texttt{NOX\_ENABLE\_TIER\_BOOST=1} opt-in for callers who want the
legacy behavior. The default reflects the calibration this evaluation
establishes; the opt-in preserves backward compatibility for downstream
pipelines that depend on it.

\subsubsection{\texorpdfstring{5.5 Claim 4 --- \texttt{source\_type}
backfill
recovery}{5.5 Claim 4 --- source\_type backfill recovery}}\label{claim-4-source_type-backfill-recovery}

Pre-backfill, \textbf{67,949 chunks (98.48\% of the corpus)} carried
\texttt{source\_type\ =\ NULL}, rendering the
\texttt{SOURCE\_TYPE\_BOOST} map inert at search time regardless of the
configured multipliers. PR \#151 backfills 11 canonical keys
(\texttt{entity}, \texttt{lesson}, \texttt{skill}, \texttt{project-doc},
\texttt{command}, \texttt{legal-template}, \texttt{personal-doc},
\texttt{session}, \texttt{note}, \texttt{external}, \texttt{other},
\texttt{ocr-cache}) by classifying each chunk via deterministic
path/prefix rules under \texttt{withOpAudit()} (audit\_id = 118),
preserving the 1,046 chunks already marked \texttt{external} (1.52\%).
The post-backfill distribution skews to \texttt{personal-doc} (32.74\%),
\texttt{skill} (19.89\%), and \texttt{session} (16.95\%), reflecting the
lived shape of the operational corpus.

In the G5 V3 matrix, A5 (source\_type only) and A10 (full minus
source\_type) both score 0.6237 --- identical to A8 --- confirming the
\texttt{SOURCE\_TYPE\_BOOST} map remained \textbf{inert by key
mismatch}, not by data absence. PR \#154 (merged 2026-05-20) updated the
boost map to the new keys (calibration ranging from
\texttt{entity\ =\ 2.0} for high-curation chunks down to
\texttt{ocr-cache\ =\ 0.7} for low-signal scanned material).

The G8 ablation (2026-05-20, PR \#177) re-ingested an isolated
\texttt{entity-eval-v2.db} (500 chunks) with source\_type values
deterministically remapped to prod-consistent vocabulary
(\texttt{entity\_file\ →\ entity}, \texttt{event\_log\ →\ lesson},
\texttt{session\_summary\ →\ session}), achieving 100\% key-match with
the SOURCE\_TYPE\_BOOST map. Results:

{\def\LTcaptype{none} % do not increment counter
\begin{longtable}[]{@{}
  >{\raggedright\arraybackslash}p{(\linewidth - 6\tabcolsep) * \real{0.2500}}
  >{\raggedright\arraybackslash}p{(\linewidth - 6\tabcolsep) * \real{0.2500}}
  >{\raggedright\arraybackslash}p{(\linewidth - 6\tabcolsep) * \real{0.2500}}
  >{\raggedright\arraybackslash}p{(\linewidth - 6\tabcolsep) * \real{0.2500}}@{}}
\toprule\noalign{}
\begin{minipage}[b]{\linewidth}\raggedright
Config
\end{minipage} & \begin{minipage}[b]{\linewidth}\raggedright
nDCG@10
\end{minipage} & \begin{minipage}[b]{\linewidth}\raggedright
Δ vs A0
\end{minipage} & \begin{minipage}[b]{\linewidth}\raggedright
Verdict
\end{minipage} \\
\midrule\noalign{}
\endhead
\bottomrule\noalign{}
\endlastfoot
A0 (no boosts) & 0.4816 & baseline & --- \\
\textbf{A5 (source\_type only)} & \textbf{0.4944} & \textbf{+2.66\%} &
LIVE --- boost contributes when keys match \\
A8 (full canonical) & 0.5798 & +20.4\% & full stack \\
A10 (full minus source\_type) & 0.5845 & +21.4\% & source\_type
\emph{removed} from full stack \\
\end{longtable}
}

A5 \textgreater{} A0 by +2.66\% empirically validates that
\texttt{SOURCE\_TYPE\_BOOST} contributes to ranking when source\_type
values match the map keys. However, A8 \textless{} A10 by −0.81\%
reveals \textbf{redundant double-boost} when \texttt{section\_boost}
(e.g., \texttt{compiled\ =\ 2.0} for entity-file truth sections) and
\texttt{SOURCE\_TYPE\_BOOST} (e.g., \texttt{entity\ =\ 2.0} for the same
chunks) stack on identical chunks, over-promoting at the cost of top-K
diversity. Per-category, the redundancy manifests as a −3.5 pp
regression on open-domain queries and a +1.4 pp gain on multi-hop.

The \textbf{G9 ablation} (2026-05-20, against the prod-flavored
\texttt{g5.db} 68k corpus, PR planned) \textbf{reproduces and amplifies}
both findings --- at production scale the boost contribution and the
redundancy are both \textbf{5× larger in magnitude} than in the
synthetic G8 set:

{\def\LTcaptype{none} % do not increment counter
\begin{longtable}[]{@{}
  >{\raggedright\arraybackslash}p{(\linewidth - 6\tabcolsep) * \real{0.2500}}
  >{\raggedright\arraybackslash}p{(\linewidth - 6\tabcolsep) * \real{0.2500}}
  >{\raggedright\arraybackslash}p{(\linewidth - 6\tabcolsep) * \real{0.2500}}
  >{\raggedright\arraybackslash}p{(\linewidth - 6\tabcolsep) * \real{0.2500}}@{}}
\toprule\noalign{}
\begin{minipage}[b]{\linewidth}\raggedright
Config
\end{minipage} & \begin{minipage}[b]{\linewidth}\raggedright
G8 (n=500)
\end{minipage} & \begin{minipage}[b]{\linewidth}\raggedright
G9 (n=68,995)
\end{minipage} & \begin{minipage}[b]{\linewidth}\raggedright
Δ G9 magnitude vs G8
\end{minipage} \\
\midrule\noalign{}
\endhead
\bottomrule\noalign{}
\endlastfoot
A0 (no boosts) & 0.4816 & 0.4108 & smaller baseline (prod diversity) \\
\textbf{A5 (source\_type only)} & \textbf{+2.66\% vs A0} &
\textbf{+14.2\% vs A0} & \textbf{5× larger} \\
A8 vs A10 (redundancy) & \textbf{−0.81\%} & \textbf{−2.6\%} & \textbf{3×
larger} \\
\end{longtable}
}

The G9 data \textbf{structurally validates} the resolution path of
mutual-exclusion logic (PR \#182, merged 2026-05-20): when a chunk has
\texttt{section\ ∈\ \{compiled,\ frontmatter,\ timeline\}} populated
(entity-file structural metadata), the \texttt{source\_type\_boost} is
gated to \texttt{0} to prevent stacking on top of
\texttt{section\_boost}. The mutex is rollback-gated via
\texttt{NOX\_DISABLE\_MUTEX\_SECTION\_SOURCE\_TYPE=1}.

The \textbf{G10 ablation} (2026-05-20, against \texttt{g9.db} 69,495
chunks) validates the mutex in production conditions:

{\def\LTcaptype{none} % do not increment counter
\begin{longtable}[]{@{}llll@{}}
\toprule\noalign{}
Config & nDCG@10 & MRR & R@10 \\
\midrule\noalign{}
\endhead
\bottomrule\noalign{}
\endlastfoot
A8' (mutex active, default) & \textbf{0.5478} & \textbf{0.5967} &
0.6183 \\
A8 (mutex disabled, rollback flag) & 0.5435 & 0.5813 & 0.6333 \\
\textbf{Δ mutex effect} & \textbf{+0.79\%} & \textbf{+2.65\%} &
−2.4\% \\
\end{longtable}
}

The mutex recovers \textasciitilde46\% of the A10 − A8 gap (where A10
fully removes \texttt{source\_type\_boost}) without removing the signal
entirely. The per-metric pattern --- MRR ↑ (top-1 quality) and R@10 ↓
(diversity) --- surfaces a deliberate trade-off: the mutex improves
precision at the cost of some recall breadth, with net positive on the
weighted nDCG@10.

The \textbf{G10b per-category breakdown} (2026-05-21, same DB, n = 100
across 5 categories) reveals which query types absorb the trade-off:

{\def\LTcaptype{none} % do not increment counter
\begin{longtable}[]{@{}llllll@{}}
\toprule\noalign{}
Category & n & nDCG@10 Δ\% & MRR Δ\% & R@10 Δ\% & Verdict \\
\midrule\noalign{}
\endhead
\bottomrule\noalign{}
\endlastfoot
single-hop & 20 & \textbf{+8.22\%} & \textbf{+13.20\%} & 0\% & strong
win \\
open-domain & 20 & \textbf{+2.42\%} & \textbf{+5.56\%} & 0\% & win \\
multi-hop & 20 & −3.95\% & −2.70\% & \textbf{−6.02\%} & regression \\
adversarial & 20 & −2.95\% & −5.88\% & 0\% & regression \\
temporal & 20 & n/a & n/a & n/a & degenerate corpus gap \\
\end{longtable}
}

The aggregate Δ (+0.43\% nDCG, +0.82\% MRR) is consistent in direction
with the G10 measurement (+0.79\% / +2.65\%) but attenuated in magnitude
--- within harness noise at n = 100, suggesting the deploy-time figure
was on the upper end of a noisy distribution. Substantively: the mutex
is a \textbf{single-hop optimizer with open-domain side benefits},
balanced against a multi-hop chain-traversal regression of −6.02\% R@10.
Net retrieval value (+0.0616 nDCG abs gain) exceeds losses (−0.0498 nDCG
abs), so the mutex stays deployed; the multi-hop regression is
documented as a follow-up candidate for \textbf{conditional mutex}
(active only when \texttt{query\_entities\ ≤\ 1}).

The \textbf{G11 trim ablation} (2026-05-20, same DB) tested whether
trimming the top SOURCE\_TYPE\_BOOST values
(\texttt{entity:\ 2.0\ →\ 1.3}, \texttt{lesson:\ 1.8\ →\ 1.2}) could
provide additive benefit over the mutex:

{\def\LTcaptype{none} % do not increment counter
\begin{longtable}[]{@{}lll@{}}
\toprule\noalign{}
Config & nDCG@10 & MRR \\
\midrule\noalign{}
\endhead
\bottomrule\noalign{}
\endlastfoot
Canonical (entity=2.0, lesson=1.8) & \textbf{0.5376} &
\textbf{0.5843} \\
Trim (entity=1.3, lesson=1.2) & 0.5337 & 0.5751 \\
Δ & \textbf{−0.73\%} & \textbf{−1.58\%} \\
\end{longtable}
}

Trim is \textbf{rejected}. The mutex already zeroes
\texttt{sourceTypeDelta} where redundancy occurs (chunks with both
\texttt{section} and high source-type); the \texttt{entity\ =\ 2.0}
still fires legitimately on legacy non-compiled chunks where the mutex
does not trigger, and trimming kills that residual signal. The
single-hop category suffered worst (−4.62\% nDCG, −7.40\% MRR),
confirming the mutex resolves redundancy \textbf{precisely}, while a
global trim over-corrects. The boost stack settles at the canonical
configuration:
\texttt{section\_boost\ ×\ source\_type\_boost\ (mutex-gated)\ ×\ additive\ salience\ v2}.

The \textbf{G10c per-style breakdown} (2026-05-21, same DB, n = 100
across 2 styles --- the dataset distinguishes \texttt{keyword} vs
\texttt{natural-language} rather than the paraphrase/literal axis the
spec anticipated) cuts the same data along a different dimension to test
whether mutex behavior is style-conditional:

{\def\LTcaptype{none} % do not increment counter
\begin{longtable}[]{@{}
  >{\raggedright\arraybackslash}p{(\linewidth - 10\tabcolsep) * \real{0.1667}}
  >{\raggedright\arraybackslash}p{(\linewidth - 10\tabcolsep) * \real{0.1667}}
  >{\raggedright\arraybackslash}p{(\linewidth - 10\tabcolsep) * \real{0.1667}}
  >{\raggedright\arraybackslash}p{(\linewidth - 10\tabcolsep) * \real{0.1667}}
  >{\raggedright\arraybackslash}p{(\linewidth - 10\tabcolsep) * \real{0.1667}}
  >{\raggedright\arraybackslash}p{(\linewidth - 10\tabcolsep) * \real{0.1667}}@{}}
\toprule\noalign{}
\begin{minipage}[b]{\linewidth}\raggedright
Style
\end{minipage} & \begin{minipage}[b]{\linewidth}\raggedright
n
\end{minipage} & \begin{minipage}[b]{\linewidth}\raggedright
nDCG@10 Δ\%
\end{minipage} & \begin{minipage}[b]{\linewidth}\raggedright
MRR Δ\%
\end{minipage} & \begin{minipage}[b]{\linewidth}\raggedright
R@10 Δ\%
\end{minipage} & \begin{minipage}[b]{\linewidth}\raggedright
Verdict
\end{minipage} \\
\midrule\noalign{}
\endhead
\bottomrule\noalign{}
\endlastfoot
natural-language & 50 & \textbf{+1.56\%} & \textbf{+3.86\%} & −1.62\% &
mutex helps \\
keyword & 50 & −0.72\% & −2.27\% & −1.06\% & mutex slightly hurts \\
\end{longtable}
}

The aggregate effect (+0.43\% nDCG, +0.82\% MRR --- identical to G10b
because G10c reuses the same A8 active vs A8 disabled detail JSONs and
re-buckets) is entirely carried by the natural-language subset. The
keyword bucket is a small drag, within the noise floor. Two cross-cuts
(style × category) surface as notable outliers: NL × single-hop
(+13.83\% nDCG, +21.32\% MRR --- the biggest individual win in the data)
and keyword × adversarial (−5.35\% nDCG, −10.0\% MRR --- the only delta
crossing the 5\% regression threshold, n = 10). Multi-hop suffers ≈ −4\%
across both styles, confirming the regression is \textbf{style-agnostic}
and motivating the conditional-mutex follow-up rather than a
style-specific routing.

Triangulated across G10 (deploy figure +0.79\% / +2.65\%), G10b
(per-category +0.43\% / +0.82\%), and G10c (per-style +0.43\% /
+0.82\%), the mutex effect is consistent in direction and on the lower
end of the original deploy measurement in magnitude --- the deploy
figure sat on the upper tail of a noisy distribution rather than
reflecting a structural shift. The architectural conclusion stands:
\textbf{keep the mutex deployed at the per-chunk level; address the
multi-hop chain-traversal regression via a conditional gate keyed on
\texttt{query\_entities}} --- executed in G10d below.

The \textbf{G10d ablation} (2026-05-21, same \texttt{g9.db} corpus, 69
495 chunks, 15 612 \texttt{kg\_entities}) tests a conditional variant of
the Hard Mutex: the per-chunk gate is suppressed when the incoming query
matches two or more entities in the KG, preserving the full boost stack
for multi-entity chain-traversal queries while keeping the mutex active
for single-entity and entity-free queries. The mechanism relies on
\texttt{query-entity-count.ts}, which performs a greedy longest-match
scan over \texttt{kg\_entities.name} at query time; the count drives
\texttt{NOX\_MUTEX\_QUERY\_ENTITY\_THRESHOLD} in \texttt{search.ts}.
Four configurations were run against the same 100-query golden set (n =
100, 5 categories × 2 styles × 10), each on an isolated endpoint (port
18803) with salience active, \textasciitilde13 min total VPS time:

{\def\LTcaptype{none} % do not increment counter
\begin{longtable}[]{@{}
  >{\raggedright\arraybackslash}p{(\linewidth - 12\tabcolsep) * \real{0.1154}}
  >{\raggedright\arraybackslash}p{(\linewidth - 12\tabcolsep) * \real{0.1154}}
  >{\raggedleft\arraybackslash}p{(\linewidth - 12\tabcolsep) * \real{0.1538}}
  >{\raggedleft\arraybackslash}p{(\linewidth - 12\tabcolsep) * \real{0.1538}}
  >{\raggedleft\arraybackslash}p{(\linewidth - 12\tabcolsep) * \real{0.1538}}
  >{\raggedleft\arraybackslash}p{(\linewidth - 12\tabcolsep) * \real{0.1538}}
  >{\raggedleft\arraybackslash}p{(\linewidth - 12\tabcolsep) * \real{0.1538}}@{}}
\toprule\noalign{}
\begin{minipage}[b]{\linewidth}\raggedright
Config
\end{minipage} & \begin{minipage}[b]{\linewidth}\raggedright
Description
\end{minipage} & \begin{minipage}[b]{\linewidth}\raggedleft
nDCG@10
\end{minipage} & \begin{minipage}[b]{\linewidth}\raggedleft
MRR
\end{minipage} & \begin{minipage}[b]{\linewidth}\raggedleft
R@10
\end{minipage} & \begin{minipage}[b]{\linewidth}\raggedleft
Δ\%nDCG vs A8'
\end{minipage} & \begin{minipage}[b]{\linewidth}\raggedleft
Δ\%MRR vs A8'
\end{minipage} \\
\midrule\noalign{}
\endhead
\bottomrule\noalign{}
\endlastfoot
A8' & G10 Hard Mutex always-on (prod baseline) & 0.5502 & 0.5992 &
0.6183 & --- & --- \\
A8d-1 & Conditional, threshold=1 & 0.5467 & 0.5856 & 0.6333 & −0.64\% &
−2.27\% \\
\textbf{A8d-2} & \textbf{Conditional, threshold=2} & \textbf{0.5577} &
\textbf{0.6074} & \textbf{0.6233} & \textbf{+1.35\%} &
\textbf{+1.37\%} \\
A8 off & Mutex fully disabled (control) & 0.5438 & 0.5806 & 0.6333 &
−1.17\% & −3.10\% \\
\end{longtable}
}

The headline result is \textbf{A8d-2 (threshold=2) wins on all three
primary metrics} over the A8' G10 baseline. A8d-1 regresses on nDCG and
MRR, pointing to a critical entity-density effect: with 15 612 entities
in \texttt{kg\_entities} --- 40× the 402 initially estimated at design
time --- a threshold of 1 is effectively always met, collapsing the
conditional back to the constant-off regime (no mutex). Threshold=2 is
the minimum noise filter that restores actual conditionality at current
entity density.

The per-category breakdown reveals the mechanism of recovery:

{\def\LTcaptype{none} % do not increment counter
\begin{longtable}[]{@{}
  >{\raggedright\arraybackslash}p{(\linewidth - 10\tabcolsep) * \real{0.1429}}
  >{\raggedright\arraybackslash}p{(\linewidth - 10\tabcolsep) * \real{0.1429}}
  >{\raggedleft\arraybackslash}p{(\linewidth - 10\tabcolsep) * \real{0.1905}}
  >{\raggedleft\arraybackslash}p{(\linewidth - 10\tabcolsep) * \real{0.1905}}
  >{\raggedleft\arraybackslash}p{(\linewidth - 10\tabcolsep) * \real{0.1905}}
  >{\raggedright\arraybackslash}p{(\linewidth - 10\tabcolsep) * \real{0.1429}}@{}}
\toprule\noalign{}
\begin{minipage}[b]{\linewidth}\raggedright
Category
\end{minipage} & \begin{minipage}[b]{\linewidth}\raggedright
n
\end{minipage} & \begin{minipage}[b]{\linewidth}\raggedleft
nDCG@10 Δ\% vs A8'
\end{minipage} & \begin{minipage}[b]{\linewidth}\raggedleft
MRR Δ\% vs A8'
\end{minipage} & \begin{minipage}[b]{\linewidth}\raggedleft
R@10 Δ\% vs A8'
\end{minipage} & \begin{minipage}[b]{\linewidth}\raggedright
Verdict
\end{minipage} \\
\midrule\noalign{}
\endhead
\bottomrule\noalign{}
\endlastfoot
multi-hop & 20 & \textbf{+1.58\%} & 0.00\% & \textbf{+3.75\%} & recovery
--- chain signal preserved \\
adversarial & 20 & \textbf{+3.04\%} & \textbf{+6.25\%} & 0.00\% &
recovery --- distractor suppression improves \\
open-domain & 20 & \textbf{+2.92\%} & \textbf{+1.59\%} & 0.00\% &
extends G10b win \\
single-hop & 20 & −3.26\% & −4.43\% & 0.00\% & trade-off \\
temporal & 20 & n/a & n/a & n/a & degenerate corpus gap (unchanged) \\
\end{longtable}
}

Multi-hop recovers because queries naming two or more entities (e.g.,
entity + associated event or related entity) now reach top-K with the
full \texttt{section\_boost\ ×\ source\_type\_boost} stack active,
enabling intermediate chain chunks to surface. Adversarial recovery is
the strongest signal: adversarial queries in the golden set tend to
mention three or more entity names as distractors, pushing
\texttt{query\_entity\_count\ ≥\ 2} and gating the mutex; the full boost
stack then differentiates gold from distractors more effectively than
the mutex-flattened ranking. The single-hop trade-off is real --- A8d-2
nDCG drops −3.26\% vs A8' --- but is bounded: single-hop performance
against the pre-mutex baseline (G10b mutex\_disabled) is still
\textbf{+3.31\% nDCG / +7.78\% MRR} (absolute: 0.5470 vs 0.5295
disabled). The conditional layer trades peak single-hop precision for
materially better worst-case behavior across multi-hop and adversarial
categories.

Latency is unaffected: P95 spread across all four configs is 2558--2573
ms (0.6\% variance), consistent with the \texttt{query-entity-count} hot
path operating at sub-millisecond cost when the entity index is warmed.

\textbf{Decision D51 verdict: ACTIVE-T2.} A8d-2 meets 6 of 8 evaluated
criteria (aggregate nDCG/MRR, multi-hop nDCG/R@10, open-domain nDCG,
adversarial nDCG). Single-hop nDCG and MRR are the two fails --- both
against the A8' baseline that represents the maximal single-hop state
--- and both remain strictly positive against the pre-mutex baseline.
The aggregate net of +1.35\% nDCG / +1.37\% MRR over the G10 baseline
justifies accepting the single-hop dilution. The canonical boost stack
in production now reads:
\texttt{section\_boost\ ×\ source\_type\_boost\ (Hard\ Mutex\ gated\ by\ query\_entity\_count\ ≤\ 2)\ ×\ salience\ v2\ additive}.

The G10d conditional gate was deployed to production on 2026-05-21 via
systemd environment drop-in
(\texttt{NOX\_MUTEX\_QUERY\_ENTITY\_THRESHOLD=2}). A smoke test across
three query archetypes confirmed correct behavior: a single-entity query
applied the mutex as expected; a multi-entity query (count ≥ 2) returned
an \texttt{entity::compiled} chunk at rank 1, confirming the mutex was
suppressed and the full boost stack served the chain; a no-entity query
bypassed the mutex entirely. Zero errors were recorded in
\texttt{journalctl} post-restart. Three rollback paths are documented
--- disabling only the conditional layer (preserving G10 hard mutex),
disabling the entire mutex, and removing the drop-in --- each executable
in under five minutes.

Triangulated across G10 (+0.79\% nDCG deploy measurement), G10b
(per-category breakdown, aggregate +0.43\%), G10c (per-style breakdown,
aggregate +0.43\%), and G10d (conditional gate, aggregate +1.35\%), the
mutex evolution follows a consistent trajectory: the per-chunk hard
mutex provided a net positive but introduced multi-hop and adversarial
regressions; the conditional layer recovers those regressions at the
cost of moderate single-hop dilution, with the aggregate strictly
improving at each step. The final deployed configuration is the most
balanced across query-category diversity the series has measured.

\subsubsection{5.6 Production deployment and
observability}\label{production-deployment-and-observability}

The G10d conditional Hard Mutex with
\texttt{NOX\_MUTEX\_QUERY\_ENTITY\_THRESHOLD=2} é a configuração
canônica em produção desde 2026-05-21. O drop-in está em
\texttt{/etc/systemd/system/nox-mem-api.service.d/override.conf}, e três
rollback paths permanecem documentados --- desabilitar apenas a camada
condicional (preservando o G10 hard mutex), desabilitar o mutex inteiro
via \texttt{NOX\_DISABLE\_MUTEX\_SECTION\_SOURCE\_TYPE=1}, ou remover o
drop-in --- cada um executável em menos de cinco minutos. O modo
\texttt{NOX\_SALIENCE\_MODE=active} (formulação aditiva v2) também está
deployado em produção, consistente com a Claim 1 da §5.2; o modo
\texttt{shadow} permanece disponível como fallback para A/B comparisons,
mas o canonical runtime usa \texttt{active}.

A camada de observabilidade F10 (Foundation observability dashboard,
decisão D53, 2026-05-21) acompanha os dois deploys em produção.
\textbf{Phase A} (\texttt{/observability/health.html}) expõe três
endpoints --- \texttt{/api/observability/health},
\texttt{/api/observability/recent-ops},
\texttt{/api/observability/canary-tail} --- com polling de 30s sobre
status do serviço, últimas operações destrutivas registradas em
\texttt{ops\_audit} (status enum
\texttt{started\ \textbar{}\ success\ \textbar{}\ failed\ \textbar{}\ crashed}),
e o tail das execuções do cron de canary. \textbf{Phase B}
(\texttt{/observability/evals.html}) consome
\texttt{/api/observability/evals} lendo \texttt{audits/data-G*/},
renderizando line charts com Chart.js sobre as séries G3 → G4 → G5 V3 →
G8 → G9 → G10 → G10b → G10c → G10d com gate annotations (D43 threshold
≥+15\% nDCG@10, D48 close, D51 verdict ACTIVE-T2). Ambas as fases
passaram smoke tests no deploy (6/6 e 5/5 respectivamente) e estão
acessíveis via Tailscale tunnel; o stack permanece lean (vanilla JS +
Chart.js CDN, sem Prometheus/Grafana/time-series DB adicional). A
leitura é em tempo real sobre o \texttt{nox-mem.db} canônico ---
qualquer regressão pós-deploy aparece nos charts dentro do próximo ciclo
de polling.

\subsubsection{5.7 Honest
characterization}\label{honest-characterization}

The +78.8\% headline is decisive for the ``Pain-weighted hybrid memory''
framing in the sense that pain is one of four additive salience
components and the full additive formulation outperforms the legacy
multiplicative one. It is \textbf{not} decisive for ``pain as a
standalone retrieval signal in hybrid mode'' --- that question is
addressed in the companion arXiv draft
(\texttt{paper/publication/paper-draft-sec4-7.md} §5.5, E10 pain
ablation: directional but not significant, Δ = +0.0065, 95\% CI
{[}−0.0143, +0.0338{]}, n = 31 on the prior R01c-v1.1 corpus). The Wave
A measurement validates the architectural choices around section-aware
ranking and additive salience composition; per-dimension causal
attribution of pain alone awaits the post-PR-\#154 ablation generation
and a corpus with broader pain distribution than the current 90.67\%
default.

A G10d evolution further refines the architectural conclusion: the
canonical boost stack
\texttt{section\_boost\ ×\ source\_type\_boost\ (Hard\ Mutex\ gated\ by\ query\_entity\_count\ ≤\ 2)\ ×\ salience\ v2\ additive}
deployed em 2026-05-21 trata a redundância identificada em G8/G9 sem
zerar o sinal completo, e recupera regressões multi-hop (+1.58\% nDCG,
+3.75\% R@10) e adversarial (+3.04\% nDCG, +6.25\% MRR) ao custo de uma
diluição contida em single-hop. A trajetória G3 → G4 → G5 V3 → G8 → G9 →
G10 → G10b → G10c → G10d demonstra disciplina de ablation: cada
generation isolou um componente ou condição, e cada decision (D43 gate,
D48 saga close, D51 ACTIVE-T2) está triangulada por código (PRs
\#150/\#151/\#153/\#154/\#177/\#181/\#182/\#198), audits
(\texttt{audits/data-G*/}), e a camada F10 que torna o resultado
verificável a qualquer momento em produção.

\begin{center}\rule{0.5\linewidth}{0.5pt}\end{center}

\subsection{6. Q4 COMPARISON --- Cross-System Benchmarking
(Pre-registered)}\label{q4-comparison-cross-system-benchmarking-pre-registered}

\begin{quote}
\textbf{Status (atualizado 2026-05-24 18h BRT):} Pre-registered skeleton
populado com o \textbf{smoke de 20 queries do Sat 2026-05-24} sobre
eval-isolated DB (5.882 chunks LoCoMo + 940 chunks LongMemEval, k=10),
validando a metodologia + confirmando que o pipeline nox-mem funciona
end-to-end no eval isolado. \textbf{Partial cross-system data adicionado
Sat 2026-05-24:} mem0 smoke completo (n=20, 500-chunk corpus cap por
cost-control, \$0.10 ingest cost estimado). Tabela §6.3 agora contém 2
linhas reais (nox-mem + mem0). Os 4 adapters restantes (Zep, Letta,
agentmemory, EverMind-AI) permanecem
\texttt{{[}PENDING\ canonical\ run{]}} até que o run canônico lande.
Princípios (§6.5), anti-cherry-pick (§6.6) e pre-registration (§6.7) são
imutáveis post-run conforme §6.7. Ref:
\texttt{{[}{[}q4-partial-cross-system-sat-2026-05-24{]}{]}}.
\end{quote}

\subsubsection{6.1 Methodology summary}\label{methodology-summary}

A §6 cobre a comparação cross-system entre nox-mem e cinco sistemas
competidores de memória persistente para agentes de IA. O execution plan
completo está documentado em
\texttt{specs/2026-05-23-Q4-comparison-execution-plan.md}
(pre-registered 2026-05-23, antes do run de Sat 2026-05-24). Os
princípios de comparação (§6.5), as garantias anti-cherry-pick (§6.6), e
a pre-registration formal (§6.7) são cravados nesta seção antes da
execução; somente as tabelas de §6.2/§6.3/§6.4 recebem números após o
run. O objetivo é satisfazer o gate D43 (\texttt{docs/DECISIONS.md}) ---
nox-mem em top-3 em ≥2 das 4 métricas chave (nDCG@10, R@10, MRR,
latência) --- destravando a GTM Phase 2.

\subsubsection{6.2 Competitors}\label{competitors}

A escolha dos cinco competidores prioriza stars no GitHub, atividade
recente de commits e overlap funcional com o escopo do nox-mem. Versões
são cravadas pré-execução para reprodutibilidade.

{\def\LTcaptype{none} % do not increment counter
\begin{longtable}[]{@{}
  >{\raggedright\arraybackslash}p{(\linewidth - 8\tabcolsep) * \real{0.2000}}
  >{\raggedright\arraybackslash}p{(\linewidth - 8\tabcolsep) * \real{0.2000}}
  >{\raggedright\arraybackslash}p{(\linewidth - 8\tabcolsep) * \real{0.2000}}
  >{\raggedright\arraybackslash}p{(\linewidth - 8\tabcolsep) * \real{0.2000}}
  >{\raggedright\arraybackslash}p{(\linewidth - 8\tabcolsep) * \real{0.2000}}@{}}
\toprule\noalign{}
\begin{minipage}[b]{\linewidth}\raggedright
System
\end{minipage} & \begin{minipage}[b]{\linewidth}\raggedright
Repo
\end{minipage} & \begin{minipage}[b]{\linewidth}\raggedright
Install path
\end{minipage} & \begin{minipage}[b]{\linewidth}\raggedright
Version pinned
\end{minipage} & \begin{minipage}[b]{\linewidth}\raggedright
Default config
\end{minipage} \\
\midrule\noalign{}
\endhead
\bottomrule\noalign{}
\endlastfoot
Mem0 & \texttt{mem0ai/mem0} & \texttt{pip\ install\ mem0ai} &
\texttt{{[}PENDING\ canonical\ run\ —\ adapter\ under\ setup{]}} &
OpenAI embeddings + Chroma vector store \\
Zep & \texttt{getzep/zep} & Docker compose (zep + postgres) &
\texttt{{[}PENDING\ canonical\ run\ —\ adapter\ under\ setup{]}} & Local
self-host mode \\
Letta (ex-MemGPT) & \texttt{letta-ai/letta} &
\texttt{pip\ install\ letta} &
\texttt{{[}PENDING\ canonical\ run\ —\ adapter\ under\ setup{]}} &
SQLite backend \\
agentmemory & \texttt{rohitg00/agentmemory} & iii-engine runtime &
\texttt{{[}PENDING\ canonical\ run\ —\ adapter\ under\ setup{]}} &
Stack-bridge mode \\
EverMind-AI & EverOS published bench & repo clone &
\texttt{{[}PENDING\ canonical\ run\ —\ adapter\ under\ setup{]}} &
Native CLI \\
\end{longtable}
}

Cada sistema roda com sua configuração default publicável (princípio
§6.5.3): nenhum competidor é tunado adversarialmente.

\subsubsection{6.3 Per-system per-dataset
results}\label{per-system-per-dataset-results}

Tabela canônica cross-system × cross-dataset. K cutoff fixado em 10 em
todos os sistemas; latência medida externamente (wall clock around
adapter call); custo derivado dos logs por-sistema (API calls × pricing
publicado).

\textbf{Sat 2026-05-24 partial cross-system smoke (20 queries combined,
dry-run-sample, eval-isolated DB).} nox-mem: full corpus (6.822 chunks =
5.882 LoCoMo + 940 LongMemEval). mem0: \textbf{500-chunk corpus cap} por
cost-control (\$0.10 ingest cost estimado; \textasciitilde8\% do corpus
completo). Caveat crítico: os números do mem0 refletem um corpus
significativamente menor --- interpretação no parágrafo abaixo.

{\def\LTcaptype{none} % do not increment counter
\begin{longtable}[]{@{}
  >{\raggedright\arraybackslash}p{(\linewidth - 16\tabcolsep) * \real{0.0857}}
  >{\raggedleft\arraybackslash}p{(\linewidth - 16\tabcolsep) * \real{0.1143}}
  >{\raggedleft\arraybackslash}p{(\linewidth - 16\tabcolsep) * \real{0.1143}}
  >{\raggedleft\arraybackslash}p{(\linewidth - 16\tabcolsep) * \real{0.1143}}
  >{\raggedleft\arraybackslash}p{(\linewidth - 16\tabcolsep) * \real{0.1143}}
  >{\raggedleft\arraybackslash}p{(\linewidth - 16\tabcolsep) * \real{0.1143}}
  >{\raggedleft\arraybackslash}p{(\linewidth - 16\tabcolsep) * \real{0.1143}}
  >{\raggedleft\arraybackslash}p{(\linewidth - 16\tabcolsep) * \real{0.1143}}
  >{\raggedleft\arraybackslash}p{(\linewidth - 16\tabcolsep) * \real{0.1143}}@{}}
\toprule\noalign{}
\begin{minipage}[b]{\linewidth}\raggedright
System
\end{minipage} & \begin{minipage}[b]{\linewidth}\raggedleft
n
\end{minipage} & \begin{minipage}[b]{\linewidth}\raggedleft
Corpus chunks
\end{minipage} & \begin{minipage}[b]{\linewidth}\raggedleft
nDCG@10
\end{minipage} & \begin{minipage}[b]{\linewidth}\raggedleft
R@10
\end{minipage} & \begin{minipage}[b]{\linewidth}\raggedleft
MRR
\end{minipage} & \begin{minipage}[b]{\linewidth}\raggedleft
p50 (ms)
\end{minipage} & \begin{minipage}[b]{\linewidth}\raggedleft
avg (ms)
\end{minipage} & \begin{minipage}[b]{\linewidth}\raggedleft
Gold hits
\end{minipage} \\
\midrule\noalign{}
\endhead
\bottomrule\noalign{}
\endlastfoot
\textbf{nox-mem} & 20 & 6.822 (full) & 0.6380 & 0.5417 & \textbf{0.3700}
& \textbf{8} & \textbf{9} & \textbf{13/20 (65\%)} \\
\textbf{mem0} (500-cap) & 20 & 500 (\textasciitilde8\%) &
\textbf{0.8569} & 0.2500 & 0.1167 & 273 & 288 & 3/20 (15\%) \\
\end{longtable}
}

Per-dataset gold-hit breakdown do nox-mem smoke: \textbf{LoCoMo 7/10
(70\%) · LongMemEval 6/10 (60\%)}.

\textbf{Interpretação do trade-off (honestidade obrigatória):}

Os dois sistemas exibem perfis opostos. nox-mem, com corpus completo
(6.822 chunks, ingest local zero-custo), produz \textbf{4× maior
hit-rate} (65\% vs 15\%) e \textbf{MRR 3× melhor} (0.37 vs 0.12) --- o
primeiro hit relevante chega antes em nox-mem. A latência de nox-mem é
\textbf{30× mais rápida} (8ms p50 vs 273ms p50), reflexo da busca local
vs chamadas à API mem0.

mem0, operando sobre apenas 500 chunks (\textasciitilde8\% do corpus),
exibe \textbf{nDCG@10 superior} (0.86 vs 0.64): os poucos hits que
retorna tendem a ser top-ranked, produzindo alta concentração de
relevância nas primeiras posições. Isso é um artefato de corpus window
menor --- com janela restrita, o sistema tem menos competição entre
resultados candidatos, o que infla o nDCG per-se mas mascara a cobertura
real (R@10 = 0.25 vs 0.54). Em produção com corpus completo e mesmo
custo de ingest, a relação nDCG pode inverter; o run canônico (corpus
uniforme, sem cap) será o árbitro desta hipótese.

Resumo executivo: \textbf{nox-mem ganha em cobertura (hits), velocidade
(latência), e first-hit quality (MRR). mem0 ganha em concentração de
relevância por-resultado (nDCG@10) dentro de uma janela de corpus
menor.} Corpus cap de 500 chunks para mem0 é cost-control explícito ---
\$0.10 estimado vs zero-cost local; dados iguais de corpus revertem
parcialmente o nDCG gap.

O smoke não disaggregou \texttt{nDCG@10} por dataset (combined-only) ---
desagregação canônica vem no run completo. Os números a seguir são da
execução canônica que ainda está em curso.

\textbf{LongMemEval n=100 (canonical):}

{\def\LTcaptype{none} % do not increment counter
\begin{longtable}[]{@{}
  >{\raggedright\arraybackslash}p{(\linewidth - 14\tabcolsep) * \real{0.0968}}
  >{\raggedleft\arraybackslash}p{(\linewidth - 14\tabcolsep) * \real{0.1290}}
  >{\raggedleft\arraybackslash}p{(\linewidth - 14\tabcolsep) * \real{0.1290}}
  >{\raggedleft\arraybackslash}p{(\linewidth - 14\tabcolsep) * \real{0.1290}}
  >{\raggedleft\arraybackslash}p{(\linewidth - 14\tabcolsep) * \real{0.1290}}
  >{\raggedleft\arraybackslash}p{(\linewidth - 14\tabcolsep) * \real{0.1290}}
  >{\raggedleft\arraybackslash}p{(\linewidth - 14\tabcolsep) * \real{0.1290}}
  >{\raggedleft\arraybackslash}p{(\linewidth - 14\tabcolsep) * \real{0.1290}}@{}}
\toprule\noalign{}
\begin{minipage}[b]{\linewidth}\raggedright
System
\end{minipage} & \begin{minipage}[b]{\linewidth}\raggedleft
nDCG@10
\end{minipage} & \begin{minipage}[b]{\linewidth}\raggedleft
R@10
\end{minipage} & \begin{minipage}[b]{\linewidth}\raggedleft
MRR
\end{minipage} & \begin{minipage}[b]{\linewidth}\raggedleft
p50 (ms)
\end{minipage} & \begin{minipage}[b]{\linewidth}\raggedleft
p95 (ms)
\end{minipage} & \begin{minipage}[b]{\linewidth}\raggedleft
p99 (ms)
\end{minipage} & \begin{minipage}[b]{\linewidth}\raggedleft
Cost/query (USD)
\end{minipage} \\
\midrule\noalign{}
\endhead
\bottomrule\noalign{}
\endlastfoot
\textbf{nox-mem} & \texttt{{[}PENDING\ canonical\ run{]}} &
\texttt{{[}PENDING\ canonical\ run{]}} &
\texttt{{[}PENDING\ canonical\ run{]}} &
\texttt{{[}PENDING\ canonical\ run{]}} &
\texttt{{[}PENDING\ canonical\ run{]}} &
\texttt{{[}PENDING\ canonical\ run{]}} & \textasciitilde\$0.00
(local) \\
Mem0 &
\texttt{{[}PENDING\ canonical\ run\ —\ full\ corpus,\ no\ 500-cap{]}} &
\texttt{{[}PENDING{]}} & \texttt{{[}PENDING{]}} & \texttt{{[}PENDING{]}}
& \texttt{{[}PENDING{]}} & \texttt{{[}PENDING{]}} &
\textasciitilde\$0.10+ ingest \\
Zep & \texttt{{[}PENDING\ canonical\ run\ —\ adapter\ under\ setup{]}} &
\texttt{{[}PENDING{]}} & \texttt{{[}PENDING{]}} & \texttt{{[}PENDING{]}}
& \texttt{{[}PENDING{]}} & \texttt{{[}PENDING{]}} &
\texttt{{[}PENDING{]}} \\
Letta & \texttt{{[}PENDING\ canonical\ run\ —\ adapter\ under\ setup{]}}
& \texttt{{[}PENDING{]}} & \texttt{{[}PENDING{]}} &
\texttt{{[}PENDING{]}} & \texttt{{[}PENDING{]}} & \texttt{{[}PENDING{]}}
& \texttt{{[}PENDING{]}} \\
agentmemory &
\texttt{{[}PENDING\ canonical\ run\ —\ adapter\ under\ setup{]}} &
\texttt{{[}PENDING{]}} & \texttt{{[}PENDING{]}} & \texttt{{[}PENDING{]}}
& \texttt{{[}PENDING{]}} & \texttt{{[}PENDING{]}} &
\texttt{{[}PENDING{]}} \\
EverMind-AI &
\texttt{{[}PENDING\ canonical\ run\ —\ adapter\ under\ setup{]}} &
\texttt{{[}PENDING{]}} & \texttt{{[}PENDING{]}} & \texttt{{[}PENDING{]}}
& \texttt{{[}PENDING{]}} & \texttt{{[}PENDING{]}} &
\texttt{{[}PENDING{]}} \\
\end{longtable}
}

\textbf{LoCoMo full (canonical):}

{\def\LTcaptype{none} % do not increment counter
\begin{longtable}[]{@{}
  >{\raggedright\arraybackslash}p{(\linewidth - 14\tabcolsep) * \real{0.0968}}
  >{\raggedleft\arraybackslash}p{(\linewidth - 14\tabcolsep) * \real{0.1290}}
  >{\raggedleft\arraybackslash}p{(\linewidth - 14\tabcolsep) * \real{0.1290}}
  >{\raggedleft\arraybackslash}p{(\linewidth - 14\tabcolsep) * \real{0.1290}}
  >{\raggedleft\arraybackslash}p{(\linewidth - 14\tabcolsep) * \real{0.1290}}
  >{\raggedleft\arraybackslash}p{(\linewidth - 14\tabcolsep) * \real{0.1290}}
  >{\raggedleft\arraybackslash}p{(\linewidth - 14\tabcolsep) * \real{0.1290}}
  >{\raggedleft\arraybackslash}p{(\linewidth - 14\tabcolsep) * \real{0.1290}}@{}}
\toprule\noalign{}
\begin{minipage}[b]{\linewidth}\raggedright
System
\end{minipage} & \begin{minipage}[b]{\linewidth}\raggedleft
nDCG@10
\end{minipage} & \begin{minipage}[b]{\linewidth}\raggedleft
R@10
\end{minipage} & \begin{minipage}[b]{\linewidth}\raggedleft
MRR
\end{minipage} & \begin{minipage}[b]{\linewidth}\raggedleft
p50 (ms)
\end{minipage} & \begin{minipage}[b]{\linewidth}\raggedleft
p95 (ms)
\end{minipage} & \begin{minipage}[b]{\linewidth}\raggedleft
p99 (ms)
\end{minipage} & \begin{minipage}[b]{\linewidth}\raggedleft
Cost/query (USD)
\end{minipage} \\
\midrule\noalign{}
\endhead
\bottomrule\noalign{}
\endlastfoot
\textbf{nox-mem} & \texttt{{[}PENDING\ canonical\ run{]}} &
\texttt{{[}PENDING\ canonical\ run{]}} &
\texttt{{[}PENDING\ canonical\ run{]}} &
\texttt{{[}PENDING\ canonical\ run{]}} &
\texttt{{[}PENDING\ canonical\ run{]}} &
\texttt{{[}PENDING\ canonical\ run{]}} & \textasciitilde\$0.00
(local) \\
Mem0 &
\texttt{{[}PENDING\ canonical\ run\ —\ full\ corpus,\ no\ 500-cap{]}} &
\texttt{{[}PENDING{]}} & \texttt{{[}PENDING{]}} & \texttt{{[}PENDING{]}}
& \texttt{{[}PENDING{]}} & \texttt{{[}PENDING{]}} &
\textasciitilde\$0.10+ ingest \\
Zep & \texttt{{[}PENDING\ canonical\ run\ —\ adapter\ under\ setup{]}} &
\texttt{{[}PENDING{]}} & \texttt{{[}PENDING{]}} & \texttt{{[}PENDING{]}}
& \texttt{{[}PENDING{]}} & \texttt{{[}PENDING{]}} &
\texttt{{[}PENDING{]}} \\
Letta & \texttt{{[}PENDING\ canonical\ run\ —\ adapter\ under\ setup{]}}
& \texttt{{[}PENDING{]}} & \texttt{{[}PENDING{]}} &
\texttt{{[}PENDING{]}} & \texttt{{[}PENDING{]}} & \texttt{{[}PENDING{]}}
& \texttt{{[}PENDING{]}} \\
agentmemory &
\texttt{{[}PENDING\ canonical\ run\ —\ adapter\ under\ setup{]}} &
\texttt{{[}PENDING{]}} & \texttt{{[}PENDING{]}} & \texttt{{[}PENDING{]}}
& \texttt{{[}PENDING{]}} & \texttt{{[}PENDING{]}} &
\texttt{{[}PENDING{]}} \\
EverMind-AI &
\texttt{{[}PENDING\ canonical\ run\ —\ adapter\ under\ setup{]}} &
\texttt{{[}PENDING{]}} & \texttt{{[}PENDING{]}} & \texttt{{[}PENDING{]}}
& \texttt{{[}PENDING{]}} & \texttt{{[}PENDING{]}} &
\texttt{{[}PENDING{]}} \\
\end{longtable}
}

Linhas que falharem smoke test (§6.5) são reportadas com nota explícita
\texttt{{[}FALHA:\ adapter\ setup\ gap{]}} em vez de omitidas,
consistente com §6.6. O run canônico atualiza estas tabelas em commit
dedicado quando 6/6 adapters estiverem prontos com corpus uniforme (sem
cap). Ref: \texttt{{[}{[}q4-partial-cross-system-sat-2026-05-24{]}{]}}.

\subsubsection{6.4 Per-category breakdown}\label{per-category-breakdown}

Decomposição por categoria de query do LongMemEval. nox-mem reporta as
seis categorias canônicas; competidores reportam idem onde a categoria
está presente no dataset original.

{\def\LTcaptype{none} % do not increment counter
\begin{longtable}[]{@{}
  >{\raggedright\arraybackslash}p{(\linewidth - 14\tabcolsep) * \real{0.0968}}
  >{\raggedleft\arraybackslash}p{(\linewidth - 14\tabcolsep) * \real{0.1290}}
  >{\raggedleft\arraybackslash}p{(\linewidth - 14\tabcolsep) * \real{0.1290}}
  >{\raggedleft\arraybackslash}p{(\linewidth - 14\tabcolsep) * \real{0.1290}}
  >{\raggedleft\arraybackslash}p{(\linewidth - 14\tabcolsep) * \real{0.1290}}
  >{\raggedleft\arraybackslash}p{(\linewidth - 14\tabcolsep) * \real{0.1290}}
  >{\raggedleft\arraybackslash}p{(\linewidth - 14\tabcolsep) * \real{0.1290}}
  >{\raggedleft\arraybackslash}p{(\linewidth - 14\tabcolsep) * \real{0.1290}}@{}}
\toprule\noalign{}
\begin{minipage}[b]{\linewidth}\raggedright
Category
\end{minipage} & \begin{minipage}[b]{\linewidth}\raggedleft
n
\end{minipage} & \begin{minipage}[b]{\linewidth}\raggedleft
nox-mem nDCG@10
\end{minipage} & \begin{minipage}[b]{\linewidth}\raggedleft
Mem0
\end{minipage} & \begin{minipage}[b]{\linewidth}\raggedleft
Zep
\end{minipage} & \begin{minipage}[b]{\linewidth}\raggedleft
Letta
\end{minipage} & \begin{minipage}[b]{\linewidth}\raggedleft
agentmemory
\end{minipage} & \begin{minipage}[b]{\linewidth}\raggedleft
EverMind-AI
\end{minipage} \\
\midrule\noalign{}
\endhead
\bottomrule\noalign{}
\endlastfoot
single-hop & \texttt{{[}PENDING\ canonical{]}} &
\texttt{{[}PENDING\ canonical{]}} & \texttt{{[}PENDING{]}} &
\texttt{{[}PENDING{]}} & \texttt{{[}PENDING{]}} & \texttt{{[}PENDING{]}}
& \texttt{{[}PENDING{]}} \\
multi-hop & \texttt{{[}PENDING\ canonical{]}} &
\texttt{{[}PENDING\ canonical{]}} & \texttt{{[}PENDING{]}} &
\texttt{{[}PENDING{]}} & \texttt{{[}PENDING{]}} & \texttt{{[}PENDING{]}}
& \texttt{{[}PENDING{]}} \\
temporal & \texttt{{[}PENDING\ canonical{]}} &
\texttt{{[}PENDING\ canonical{]}} & \texttt{{[}PENDING{]}} &
\texttt{{[}PENDING{]}} & \texttt{{[}PENDING{]}} & \texttt{{[}PENDING{]}}
& \texttt{{[}PENDING{]}} \\
adversarial & \texttt{{[}PENDING\ canonical{]}} &
\texttt{{[}PENDING\ canonical{]}} & \texttt{{[}PENDING{]}} &
\texttt{{[}PENDING{]}} & \texttt{{[}PENDING{]}} & \texttt{{[}PENDING{]}}
& \texttt{{[}PENDING{]}} \\
open-domain & \texttt{{[}PENDING\ canonical{]}} &
\texttt{{[}PENDING\ canonical{]}} & \texttt{{[}PENDING{]}} &
\texttt{{[}PENDING{]}} & \texttt{{[}PENDING{]}} & \texttt{{[}PENDING{]}}
& \texttt{{[}PENDING{]}} \\
numeric & \texttt{{[}PENDING\ canonical{]}} &
\texttt{{[}PENDING\ canonical{]}} & \texttt{{[}PENDING{]}} &
\texttt{{[}PENDING{]}} & \texttt{{[}PENDING{]}} & \texttt{{[}PENDING{]}}
& \texttt{{[}PENDING{]}} \\
\end{longtable}
}

O \textbf{smoke de Sat 2026-05-24} não disaggregou per-category
(combined-only sobre 20 queries), portanto §6.4 inteira aguarda o run
canônico de 100 queries × 2 datasets × 6 sistemas. Quando uma categoria
não tem queries suficientes (n \textless{} 10) em algum dataset, a
célula recebe \texttt{n/a} em vez de extrapolação, evitando o tipo de
regressão que aparece em §5.5 (temporal \texttt{n/a} na G10b por corpus
gap degenerado).

\subsubsection{6.5 Fair-comparison
principles}\label{fair-comparison-principles}

A comparação obedece a princípios padronizados pela literatura de
benchmark publicado (EverMemBench, BEIR, MTEB):

\begin{enumerate}
\def\labelenumi{\arabic{enumi}.}
\tightlist
\item
  \textbf{Corpus idêntico.} Todos os sistemas recebem o mesmo
  \texttt{chunks.text} ingerido via a API nativa de cada um. Nenhum
  sistema recebe versão ``otimizada'' do corpus.
\item
  \textbf{Eval set idêntico.} Mesmas queries, mesmos gold sets, mesmo
  random seed (\texttt{42} para shuffle do LongMemEval).
\item
  \textbf{Defaults nativos por sistema.} Cada competidor roda com sua
  configuração default publicável. Não tunamos competidores
  adversarialmente para perder; se default config é o que é publicado, é
  o que é avaliado.
\item
  \textbf{K cutoff fixo em 10.} Alguns sistemas defaultam para 5 ou 20;
  todos são forçados a \texttt{k=10} para comparabilidade.
\item
  \textbf{Embeddings provider nativo por sistema.} nox-mem usa Gemini
  3072d; cada competidor usa seu provider default. Uma variação
  \texttt{all-Gemini} é planejada como side experiment opcional,
  deferred porque o smoke de Sat 2026-05-24 consumiu o time-box antes do
  experiment
  (\texttt{{[}deferred\ —\ Sun\ 2026-05-25\ ou\ follow-up\ post-launch{]}}).
\item
  \textbf{Hardware uniforme.} Mesmo VPS (Hostinger 8 cores / 16 GB RAM),
  localhost between systems exceto chamadas a embeddings APIs externas.
\end{enumerate}

Cada adapter passa um smoke test pré-run:

\begin{Shaded}
\begin{Highlighting}[]
\NormalTok{result }\OperatorTok{=}\NormalTok{ adapter.search(}\StringTok{"test query"}\NormalTok{, k}\OperatorTok{=}\DecValTok{5}\NormalTok{)}
\ControlFlowTok{assert} \BuiltInTok{len}\NormalTok{(result) }\OperatorTok{\textgreater{}=} \DecValTok{1}
\ControlFlowTok{assert} \BuiltInTok{all}\NormalTok{(}\StringTok{\textquotesingle{}id\textquotesingle{}} \KeywordTok{in}\NormalTok{ r }\KeywordTok{and} \StringTok{\textquotesingle{}score\textquotesingle{}} \KeywordTok{in}\NormalTok{ r }\ControlFlowTok{for}\NormalTok{ r }\KeywordTok{in}\NormalTok{ result)}
\end{Highlighting}
\end{Shaded}

Adapter que falhar smoke é documentado como gap
(\texttt{{[}FALHA:\ \textless{}razão\textgreater{}{]}}) em vez de
omitido --- consistente com §6.6.

\subsubsection{6.6 Anti-cherry-pick
statement}\label{anti-cherry-pick-statement}

Para evitar viés de seleção retroativo:

\begin{itemize}
\tightlist
\item
  \textbf{Todas as 6 categorias reportadas.} Nenhuma é omitida porque o
  resultado é desfavorável.
\item
  \textbf{Ambos os datasets reportados.} LongMemEval n=100 + LoCoMo
  full, lado a lado. Não escolhemos o que beneficia.
\item
  \textbf{Latência worst-case reportada.} p50 + p95 + p99 explícitos.
  Não publicamos apenas p50.
\item
  \textbf{Per-system per-category transparency.} A tabela de §6.4 expõe
  cada combinação; não há linha agregada que mascare um padrão.
\item
  \textbf{Gaps documentados.} Sistemas que falharem setup recebem nota
  explícita; a comparação roda sem o sistema faltante, mas o gap é
  registrado em \texttt{docs/COMPARISON.md}.
\end{itemize}

\subsubsection{6.7 Pre-registration}\label{pre-registration}

A metodologia desta seção está cravada no
\texttt{specs/2026-05-23-Q4-comparison-execution-plan.md} antes do run
de Sat 2026-05-24. O \textbf{smoke de Sat 2026-05-24 15h30 BRT}
preencheu a primeira linha de §6.3 (nox-mem combined: nDCG@10=0.6380,
p50=8ms, gold-hit 13/20 em 20 queries dry-run-sample) e validou que o
pipeline de retrieval funciona end-to-end em eval-isolated DB. O
\textbf{partial cross-system smoke de Sat 2026-05-24 18h BRT} adicionou
a linha mem0 (n=20, 500-chunk corpus cap): nDCG@10=0.8569, p50=273ms,
gold-hit 3/20 (15\%) --- com interpretação explícita do trade-off
coverage vs concentração em §6.3. O \textbf{run canônico} ainda está em
curso e atualiza as linhas competidoras
\texttt{{[}PENDING\ canonical\ run{]}} em §6.3 + a totalidade de §6.4
quando os 6 adapters estiverem prontos com corpus uniforme. Princípios
(§6.5), anti-cherry-pick (§6.6) e a estrutura geral desta seção são
imutáveis post-run. Qualquer ajuste metodológico identificado durante a
execução é documentado como follow-up explícito em
\texttt{docs/COMPARISON.md} em vez de retroagido aqui. Refs:
\texttt{{[}{[}q4-smoke-sat-2026-05-24-real-numbers{]}{]}} ·
\texttt{{[}{[}q4-partial-cross-system-sat-2026-05-24{]}{]}}.

A decisão D43 (\texttt{docs/DECISIONS.md}) define o gate de aprovação:
nox-mem em top-3 em ≥2 das 4 métricas chave (nDCG@10, R@10, MRR,
latência). Atendido o gate, GTM Phase 2 está destravada conforme
\texttt{docs/ROADMAP.md} §7. Não atendido, a sessão de Sun 2026-05-25
produz um plano de remediação (ajustes pre-launch) em vez de launch
direto.

\begin{center}\rule{0.5\linewidth}{0.5pt}\end{center}

\subsection{7. Limitations and Future
Work}\label{limitations-and-future-work}

\subsubsection{7.1 Limitations}\label{limitations}

\paragraph{L1 --- Explicit-ingestion dependency (no zero-shot corpus
coverage)}\label{l1-explicit-ingestion-dependency-no-zero-shot-corpus-coverage}

nox-mem retrieves only what has been explicitly ingested via
\texttt{ingestFile()}, \texttt{ingest-entity}, or the inotifywait
watcher pipeline. There is no mechanism to answer queries over arbitrary
external corpora at query time. This is a deliberate design constraint
--- the system is optimized for an agent's \emph{own} accumulated
memory, not general-purpose retrieval augmentation --- but it means that
coverage is bounded by ingestion discipline. A corpus that has never
been ingested produces zero recall regardless of query quality. Users
bootstrapping the system must explicitly run \texttt{nox-mem\ reindex}
over existing files before the hybrid search layer is useful. Ref:
\texttt{specs/2026-03-14-nox-memory-system-design.md}, ingestion
pipeline §3.1.

\paragraph{L2 --- Gemini API dependency for embeddings (cost + outbound
network)}\label{l2-gemini-api-dependency-for-embeddings-cost-outbound-network}

The semantic retrieval layer (Layer 2) depends on Google's
\texttt{gemini-embedding-001} model (3072 dimensions). This introduces
two constraints: (a) every vectorization call requires outbound network
access and a valid \texttt{GEMINI\_API\_KEY}, meaning an air-gapped
deployment falls back to FTS5-only retrieval with no semantic recall;
(b) API cost scales with corpus size --- at the Sat 2026-05-24 corpus of
\textasciitilde69k chunks, a full re-vectorization pass takes
approximately 30--40 minutes at quota limits of the free tier. The
Autonomy pillar of the Q/A/P strategy explicitly calls out ``provider
your choice, zero vendor lock-in'' as a long-term goal; local embedding
substitution (e.g., \texttt{nomic-embed-text} via Ollama) is
architecturally feasible but not validated against the canonical eval
set. BYOK partial autonomy is available: users who supply their own
Gemini API key operate without per-query billing exposure in the default
free tier. Ref: \texttt{docs/DECISIONS.md} (model selection),
\texttt{{[}{[}default-flash-lite-for-agent-infra-tasks{]}{]}}.

\paragraph{L3 --- Single-instance architecture (no distributed sharding
or
replication)}\label{l3-single-instance-architecture-no-distributed-sharding-or-replication}

The system runs on a single SQLite file per agent database. WAL mode
provides concurrent read safety, but there is no horizontal sharding, no
replication across nodes, and no distributed coordination layer. The
current production corpus (69k chunks across 7 databases on a 4-vCPU /
8GB KVM4) operates comfortably within these bounds, but the architecture
does not generalize to multi-tenant deployments or corpora significantly
exceeding the single-node memory/storage envelope. Distributed SQLite
extensions (e.g., \texttt{cr-sqlite} CRDT-based replication) exist but
are explicitly out of scope for v1. This is a known architectural
decision, not an oversight. Ref: \texttt{docs/DECISIONS.md}
(single-instance rationale).

\paragraph{L4 --- No write-side concurrency control
(last-writer-wins)}\label{l4-no-write-side-concurrency-control-last-writer-wins}

Chunk ingestion operates under an optimistic concurrency model:
\texttt{ingestFile()} deletes existing chunks for the source file and
re-inserts in a single transaction, but there is no row-level locking or
version fence against concurrent ingest of the same source file from two
processes. In practice the inotifywait watcher and manual CLI calls
rarely overlap, and the WAL journal prevents data corruption; however,
two concurrent ingest calls on the same file produce non-deterministic
chunk counts. The \texttt{withOpAudit()} wrapper
(\texttt{src/lib/op-audit.ts}) does not add a mutual-exclusion layer for
ingest --- it targets destructive bulk operations (reindex, consolidate,
crystallize). Production mitigations are operational (systemd service
prevents concurrent watcher processes; cron stagger of 5 minutes between
agents), not architectural. Ref: \texttt{docs/INCIDENTS.md\#2026-04-25},
\texttt{{[}{[}a1-op-audit-module{]}{]}}.

\paragraph{L5 --- Evaluation sample size and canonical run gap (n=20
smoke vs n=100
target)}\label{l5-evaluation-sample-size-and-canonical-run-gap-n20-smoke-vs-n100-target}

The Sat 2026-05-24 cross-system comparison (§6.3) is based on a 20-query
smoke over an eval-isolated DB (5,882 LoCoMo + 940 LongMemEval chunks),
not the pre-registered canonical 100-query × 2-dataset × 6-system run.
The canonical run was in progress at the time of this writing (5/6
competitor adapters under setup). All competitive figures for Mem0, Zep,
Letta, agentmemory, and EverMind-AI in §6 carry
\texttt{{[}PENDING\ canonical\ run{]}} tags and should not be treated as
settled results. The nox-mem smoke figure (nDCG@10 = 0.6380 combined,
p50 = 12 ms) is validated on the methodology but not directly comparable
to the G5 V3 entity-eval figure (0.6237) because the eval corpus and
query set differ. Ref:
\texttt{specs/2026-05-23-Q4-comparison-execution-plan.md},
\texttt{{[}{[}q4-smoke-sat-2026-05-24-real-numbers{]}{]}}.

\paragraph{L6 --- Cross-system comparison is methodologically
partial}\label{l6-cross-system-comparison-is-methodologically-partial}

Three of five competitors could not be evaluated against the full
canonical corpus at the time of writing. Mem0 ran against a 500-chunk
corpus cap imposed by cost-control constraints (\$0.10 estimated ingest
cost at full corpus), producing a nDCG@10 = 0.8569 on 20 queries that
cannot be directly compared to nox-mem's full-corpus score --- a
smaller, more concentrated corpus tends to inflate nDCG for systems that
retrieve all relevant documents. Zep and Letta require Docker compose
setups that were in progress. EverMind-AI's repository was unavailable
at time of access. The ``concentration vs coverage trade-off'' (high
nDCG on capped corpus vs recall breadth on full corpus) is a genuine
open question for per-system fair comparison, not a methodological
failure. Ref: §6.3 anti-cherry-pick statement,
\texttt{docs/COMPARISON.md},
\texttt{{[}{[}q4-partial-cross-system-sat-2026-05-24{]}{]}}.

\paragraph{L7 --- Latency comparison conflates transport
classes}\label{l7-latency-comparison-conflates-transport-classes}

The Sat 2026-05-24 latency figures compare nox-mem (local FTS5 +
sqlite-vec with Gemini API call for query embedding) against competitor
adapters that go over HTTP or Python SDK to local Docker services. The
nox-mem p50 = 12 ms figure reflects localhost UNIX-domain retrieval; the
Mem0 p50 = 273 ms reflects a Docker-in-Docker HTTP call. These are not
the same transport class. A valid latency comparison requires a
normalized transport --- either all systems behind the same HTTP
gateway, or all measured at the SDK level without network hop
differences. The §6 latency figures are reported with this caveat in the
per-system notes and should not be interpreted as head-to-head speed
claims. Ref: \texttt{{[}{[}q3-latency-numbers-2026-05-18{]}{]}},
\texttt{docs/PERFORMANCE.md}.

\paragraph{L8 --- Pain signal is directional but not statistically
significant in
isolation}\label{l8-pain-signal-is-directional-but-not-statistically-significant-in-isolation}

The ``pain-weighted hybrid memory'' framing rests primarily on the
additive salience formula (§5.2) and section-aware ranking (§5.3). The
\texttt{pain} dimension contributes W\_PAIN = 0.10 of the salience
weight, but its isolated causal contribution has not been validated to
statistical significance: the E10 pain ablation
(\texttt{paper/publication/paper-draft-sec4-7.md} §5.5) reports Δ =
+0.0065 with 95\% CI {[}−0.0143, +0.0338{]} on n = 31, directional but
not significant. The current production corpus has 90.67\% of chunks at
the default \texttt{pain\ =\ 0.2}, providing insufficient variance for a
precise estimate. A definitive pain signal ablation requires a corpus
where pain scores span the full {[}0.1, 1.0{]} range. Ref: §5.7 honest
characterization, \texttt{{[}{[}d47-path-c-decision{]}{]}}.

\begin{center}\rule{0.5\linewidth}{0.5pt}\end{center}

\subsubsection{7.2 Future Work}\label{future-work}

\paragraph{F1 --- A2 Tier 3 P5: production-ready encrypted memory (in
flight)}\label{f1-a2-tier-3-p5-production-ready-encrypted-memory-in-flight}

Phase 5 of the A2 Tier 3 roadmap targets a full SQLCipher-encrypted
memory store with Ed25519-signed audit checkpoints (P4 deployed via PR
\#294). The signed checkpoint chain enables tamper-evident audit across
destructive operations (\texttt{reindex}, \texttt{consolidate},
\texttt{crystallize}) without requiring a central trust authority. P5
closes the Tier 3 arc by integrating encryption key management with the
existing \texttt{withOpAudit()} wrapper and the Tier 3 reads-audit layer
(P3, PR \#292/\#293). Target deployment: post-GTM Phase 2 launch,
estimated Sun 2026-05-25. Ref:
\texttt{specs/2026-05-24-A2-tier3-crypto-audit-RECON.md},
\texttt{docs/A2-TIER3-MIGRATION-RUNBOOK.md},
\texttt{{[}{[}a1-op-audit-module{]}{]}}.

\paragraph{F2 --- F10 Phase C/D: shadow tracker empirical A/B for
ranking
changes}\label{f2-f10-phase-cd-shadow-tracker-empirical-ab-for-ranking-changes}

F10 Phase A (\texttt{/observability/health.html}) and Phase B
(\texttt{/observability/evals.html}) are deployed (§5.6, decision D53).
Phase C targets a shadow-mode query logger that captures production
queries, executes them against a candidate ranking config in parallel,
and accumulates query-level nDCG deltas before any promotion decision.
Phase D operationalizes this into a pre-promotion gate: any ranking
change (boost weight adjustment, mutex threshold, salience weight) that
has not accumulated ≥50 shadow queries with p \textless{} 0.05
improvement is blocked from reaching the production endpoint. This
closes the observability gap identified in
\texttt{{[}{[}ship-ranking-changes-in-shadow-mode-first{]}{]}} ---
currently the shadow mode is a flag toggle, not an integrated eval
pipeline. Ref: \texttt{docs/ROADMAP.md} F10, decision D53, PR
\#207/\#212.

\paragraph{F3 --- Per-method benchmark Phase B: cross-method nDCG
optimization}\label{f3-per-method-benchmark-phase-b-cross-method-ndcg-optimization}

The Q4 per-method benchmark (§6,
\texttt{specs/2026-05-21-per-method-benchmark-comparison.md})
establishes the cross-system baseline. Phase B targets per-query-type
boost calibration: given that keyword queries respond differently from
natural-language queries (G10c §5.5), and single-hop vs multi-hop have
opposing mutex trade-offs (G10b §5.5), a routing layer that selects
ranking parameters based on query-type classification has measurable
potential upside. Estimated Lab Q1 item. Ref:
\texttt{specs/2026-05-21-per-method-benchmark-comparison.md},
\texttt{{[}{[}g10c-per-style-mutex-2026-05-21{]}{]}}.

\paragraph{F4 --- EverMemBench equivalence: honest comparison against
EverMind-AI
dataset}\label{f4-evermembench-equivalence-honest-comparison-against-evermind-ai-dataset}

\texttt{{[}{[}everos-honest-comparison-benchmark-gap{]}{]}} identifies
that EverMind-AI publishes standardized results on EverMemBench
(EverCore 83\% LongMemEval / 93\% LoCoMo; HyperMem 92.73\% LongMemEval).
Running nox-mem on EverMemBench with the same evaluation protocol closes
the benchmark gap and provides a reviewer-grade comparison for the arXiv
submission. This is gated on EverMind-AI repository availability
(currently unavailable) or an alternative comparable dataset. Estimated
Lab Q1 priority if the repository returns. Ref:
\texttt{{[}{[}everos-benchmark-publisher-competitor{]}{]}},
\texttt{{[}{[}benchmark-gap-longmemeval-locomo{]}{]}}.

\paragraph{F5 --- Neural reranker: cross-encoder rerank
post-RRF}\label{f5-neural-reranker-cross-encoder-rerank-post-rrf}

The current retrieval stack terminates at RRF fusion (§4.1). A
cross-encoder reranker --- receiving the top-K RRF candidates and the
original query as a pair --- is the standard next step in multi-stage
retrieval and typically yields +3--8\% nDCG@10 over bi-encoder baselines
(see e.g., Nogueira \& Cho 2019 on MS MARCO). The Autonomy constraint
(\texttt{{[}{[}neural-reranker-evolution-vector{]}{]}}) favors a
locally-runnable cross-encoder (e.g.,
\texttt{cross-encoder/ms-marco-MiniLM-L-6-v2} via sentence-transformers,
\textasciitilde66MB) over a cloud inference call, keeping the retrieval
stack fully offline-capable. Estimated Lab Q1/Q2. Ref:
\texttt{{[}{[}neural-reranker-as-vetor-evolutivo-pos-rrf{]}{]}},
\texttt{docs/ROADMAP.md} Lab Q1.

\paragraph{F6 --- Lab Q1 scale validation: 250k chunk
corpus}\label{f6-lab-q1-scale-validation-250k-chunk-corpus}

The current production corpus is 69,495 chunks (as of G10, 2026-05-20).
The Lab Q1 roadmap targets a 250k chunk corpus to validate: (a)
sqlite-vec ANN recall at scale (current exact-search; approximate search
becomes necessary past \textasciitilde100k vectors at reasonable latency
targets); (b) salience formula stability (the recency component decays
over a longer history window); (c) FTS5 BM25 IDF calibration (with more
documents, rare-term IDF weights shift). The
\texttt{{[}{[}lab-q1-scale-250k{]}{]}} item has no committed spec yet;
it is gated on \texttt{NOX\_SALIENCE\_MODE=active} remaining stable
through the GTM Phase 2 feedback cycle. Ref: \texttt{docs/ROADMAP.md}
Lab Q1, \texttt{{[}{[}q-a-p-pillars-strategic-pivot-2026-05-17{]}{]}}.

\paragraph{F7 --- Multilingual corpus coverage: Portuguese and
Spanish}\label{f7-multilingual-corpus-coverage-portuguese-and-spanish}

The current evaluation corpus is English-dominant (LongMemEval and
LoCoMo are English datasets; the internal entity-eval golden set mixes
English and Portuguese). The FTS5 \texttt{unicode61} tokenizer with
porter stemmer does not stem Portuguese or Spanish tokens correctly
(e.g., ``decisão'' stems to ``decisa'' rather than ``decid-''; Spanish
gerunds lose morphological overlap). The Gemini semantic layer partially
compensates via cross-lingual embedding space, but there is no explicit
multilingual evaluation. A Portuguese golden set is a natural next step
given the production operational language of the corpus. This is
deferred to post-launch community feedback intake.

\paragraph{F8 --- GTM Phase 2 launch and community feedback
intake}\label{f8-gtm-phase-2-launch-and-community-feedback-intake}

The Q/A/P roadmap (decision
\texttt{{[}{[}qap-pillars-strategic-pivot-2026-05-17{]}{]}}) defines GTM
Phase 2 as gated on D43 (nox-mem top-3 in ≥2 of 4 key metrics ---
nDCG@10, R@10, MRR, latency). The target launch date is Wed 2026-06-03,
conditional on the canonical Q4 run completing and D43 passing.
Post-launch, community feedback from the OSS release is expected to
surface real-world limitation patterns not visible in the synthetic
golden sets (e.g., corpora with high image-to-text OCR content,
multi-language mixes, or very short memory fragments \textless{} 20
words that the current chunker merges). The feedback cycle directly
informs Lab Q2 priorities. Ref: \texttt{docs/ROADMAP.md} GTM Phase 2,
\texttt{docs/gtm/},
\texttt{{[}{[}overnight-automode-push-pattern{]}{]}}.

\begin{center}\rule{0.5\linewidth}{0.5pt}\end{center}

\subsection{8. Knowledge Graph v2}\label{knowledge-graph-v2}

\subsubsection{8.1 Entity Extraction}\label{entity-extraction}

\textbf{v1 (Regex-based)}: Used hardcoded regular expressions for 3
entity types (person, project, agent) with a static alias map for name
normalization. Limited to predefined names, producing 26 entities.

\textbf{v2 (LLM-powered)}: Uses Ollama llama3.2:3b with a structured
extraction prompt. Each chunk is processed with temperature 0.1 for
deterministic output. The LLM returns JSON with entities (name + type)
and relations (source + relation + target).

Extraction results after processing 866 chunks:

{\def\LTcaptype{none} % do not increment counter
\begin{longtable}[]{@{}llll@{}}
\toprule\noalign{}
Metric & Regex v1 & LLM v2 & Improvement \\
\midrule\noalign{}
\endhead
\bottomrule\noalign{}
\endlastfoot
Entities & 26 & 384 & 14.8x \\
Relations & 59 & 529 & 9.0x \\
Entity Types & 3 & 11 & 3.7x \\
\end{longtable}
}

\textbf{Entity Type Distribution:}

{\def\LTcaptype{none} % do not increment counter
\begin{longtable}[]{@{}lll@{}}
\toprule\noalign{}
Type & Count & Description \\
\midrule\noalign{}
\endhead
\bottomrule\noalign{}
\endlastfoot
project & 109 & Software projects, products, repos \\
tool & 67 & Libraries, frameworks, CLI tools \\
concept & 54 & Abstract ideas, patterns, methodologies \\
person & 53 & Team members, contacts, stakeholders \\
organization & 50 & Companies, teams, departments \\
agent & 45 & AI agents in the fleet \\
location & 2 & Geographic references \\
other & 4 & Device, currency, date, computer \\
\end{longtable}
}

\subsubsection{8.2 Temporal Decay and TTL}\label{temporal-decay-and-ttl}

Relations have a 90-day time-to-live (TTL) from creation. The confidence
decay mechanism operates as follows:

\begin{enumerate}
\def\labelenumi{\arabic{enumi}.}
\tightlist
\item
  Relations start with confidence 0.8 (extracted) or 0.9 (confirmed)
\item
  Every 30 days without re-confirmation, confidence drops by 0.1
\item
  Relations below 0.3 confidence receive accelerated 7-day expiry
\item
  Expired relations are deleted during \texttt{kg-prune} execution
\item
  Re-confirmation (observing the same relation in new chunks) resets
  confidence to 0.9 and extends TTL by 90 days
\end{enumerate}

This mechanism ensures the knowledge graph naturally forgets stale
information while reinforcing actively observed patterns.

\subsubsection{8.3 Decision Versioning}\label{decision-versioning}

Architectural decisions are tracked with full version history in the
\texttt{decision\_versions} table. Each decision has a unique key (e.g.,
\texttt{dedup-strategy}, \texttt{fallback-chain}) and supports:

\begin{itemize}
\tightlist
\item
  Version chains with supersession tracking
\item
  Authorship attribution
\item
  Source file provenance
\item
  Current vs.~historical querying
\end{itemize}

10 decisions are currently tracked, covering API key management, LLM
fallback chains, embedding model selection, agent isolation strategy,
and synchronization schedules.

\subsubsection{8.4 Graph Traversal}\label{graph-traversal}

The \texttt{findPath()} function implements BFS (Breadth-First Search)
to discover shortest paths between any two entities. This enables
queries like ``How is Toto connected to nox-mem?'' which traverses
person → project → tool → agent relationships. Maximum depth is
configurable (default: 4 hops).

\begin{center}\rule{0.5\linewidth}{0.5pt}\end{center}

\subsection{9. Cross-Agent Intelligence}\label{cross-agent-intelligence}

\subsubsection{9.1 Agent Expertise
Profiling}\label{agent-expertise-profiling}

Each agent's memory is analyzed to determine its unique expertise based
on chunk type distribution. The dominant chunk type determines the
agent's strength category:

\begin{itemize}
\tightlist
\item
  \textbf{daily} → ``Daily operations \& activity logging''
\item
  \textbf{team} → ``Team coordination \& shared knowledge''
\item
  \textbf{decision} → ``Decision tracking \& rationale''
\item
  \textbf{lesson} → ``Lessons learned \& pattern recognition''
\end{itemize}

Profiles include chunk counts, type breakdowns, top topics (via FTS5
term frequency), and last activity dates.

\subsubsection{9.2 Knowledge Sharing}\label{knowledge-sharing}

The \texttt{pullInsightsFrom()} function enables any agent to query
lessons and decisions from other agents without direct database access.
This creates a knowledge transfer mechanism where, for example, Cipher
(Security) can learn from Forge's (Code Reviewer) past code review
decisions.

\texttt{pullAllInsights()} aggregates insights across all agents, sorted
by date, providing a fleet-wide learning feed.

\subsubsection{9.3 Cross-Agent Knowledge Graph
Merge}\label{cross-agent-knowledge-graph-merge}

\texttt{mergeCrossKnowledgeGraphs()} scans all agent databases for
kg\_entities and kg\_relations tables, merging them into a unified
entity view. Entities are matched by type + lowercase name. The output
shows which entities are known to which agents and their combined
mention counts, enabling identification of shared knowledge
vs.~agent-specific expertise.

\begin{center}\rule{0.5\linewidth}{0.5pt}\end{center}

\subsection{10. MCP Server Interface}\label{mcp-server-interface}

nox-mem exposes 14 tools via the Model Context Protocol (MCP) over stdio
(JSON-RPC 2.0):

{\def\LTcaptype{none} % do not increment counter
\begin{longtable}[]{@{}
  >{\raggedright\arraybackslash}p{(\linewidth - 4\tabcolsep) * \real{0.2069}}
  >{\raggedright\arraybackslash}p{(\linewidth - 4\tabcolsep) * \real{0.3448}}
  >{\raggedright\arraybackslash}p{(\linewidth - 4\tabcolsep) * \real{0.4483}}@{}}
\toprule\noalign{}
\begin{minipage}[b]{\linewidth}\raggedright
Tool
\end{minipage} & \begin{minipage}[b]{\linewidth}\raggedright
Category
\end{minipage} & \begin{minipage}[b]{\linewidth}\raggedright
Description
\end{minipage} \\
\midrule\noalign{}
\endhead
\bottomrule\noalign{}
\endlastfoot
nox\_mem\_search & Retrieval & Hybrid search (FTS5 + semantic + RRF) \\
nox\_mem\_stats & Monitoring & Database statistics and health \\
nox\_mem\_primer & Context & Session recovery summary
(\textasciitilde500 tokens) \\
nox\_mem\_ingest & Ingestion & Index a file into memory \\
nox\_mem\_cross\_search & Cross-Agent & Search across all 7 databases \\
nox\_mem\_cross\_stats & Cross-Agent & Chunk counts per agent \\
nox\_mem\_metrics & Monitoring & Daily observability metrics \\
nox\_mem\_kg\_build & KG & Build knowledge graph from chunks \\
nox\_mem\_kg\_query & KG & Query entity and its relations \\
nox\_mem\_kg\_stats & KG & Knowledge graph statistics \\
nox\_mem\_agent\_profiles & Intelligence & Agent expertise profiles \\
nox\_mem\_cross\_kg & Intelligence & Merged cross-agent knowledge
graph \\
nox\_mem\_kg\_path & Intelligence & BFS path between entities \\
nox\_mem\_self\_improve & Analysis & Contradiction detection, pattern
analysis \\
\end{longtable}
}

\begin{center}\rule{0.5\linewidth}{0.5pt}\end{center}

\subsection{11. HTTP API Server}\label{http-api-server}

A lightweight HTTP API (Node.js built-in \texttt{http} module, zero
dependencies) runs on port 18800, exposing memory data to the React
dashboard:

{\def\LTcaptype{none} % do not increment counter
\begin{longtable}[]{@{}
  >{\raggedright\arraybackslash}p{(\linewidth - 4\tabcolsep) * \real{0.3571}}
  >{\raggedright\arraybackslash}p{(\linewidth - 4\tabcolsep) * \real{0.2857}}
  >{\raggedright\arraybackslash}p{(\linewidth - 4\tabcolsep) * \real{0.3571}}@{}}
\toprule\noalign{}
\begin{minipage}[b]{\linewidth}\raggedright
Endpoint
\end{minipage} & \begin{minipage}[b]{\linewidth}\raggedright
Method
\end{minipage} & \begin{minipage}[b]{\linewidth}\raggedright
Response
\end{minipage} \\
\midrule\noalign{}
\endhead
\bottomrule\noalign{}
\endlastfoot
\texttt{/api/health} & GET & System health: chunks, consolidation,
vector coverage, services, KG stats, DB size \\
\texttt{/api/agents} & GET & Agent expertise profiles array \\
\texttt{/api/kg} & GET & Knowledge graph entities and relations \\
\texttt{/api/kg/path?from=X\&to=Y} & GET & BFS shortest path between
entities \\
\texttt{/api/search?q=QUERY\&limit=N} & GET & Hybrid search results \\
\texttt{/api/cross-kg} & GET & Merged cross-agent knowledge graph \\
\end{longtable}
}

CORS headers are set for cross-origin access from the Vercel-hosted
dashboard.

\begin{center}\rule{0.5\linewidth}{0.5pt}\end{center}

\subsection{12. Operational
Infrastructure}\label{operational-infrastructure}

\subsubsection{12.1 Cron Schedule}\label{cron-schedule}

24 cron jobs manage automated operations:

{\def\LTcaptype{none} % do not increment counter
\begin{longtable}[]{@{}
  >{\raggedright\arraybackslash}p{(\linewidth - 6\tabcolsep) * \real{0.1935}}
  >{\raggedright\arraybackslash}p{(\linewidth - 6\tabcolsep) * \real{0.3548}}
  >{\raggedright\arraybackslash}p{(\linewidth - 6\tabcolsep) * \real{0.1613}}
  >{\raggedright\arraybackslash}p{(\linewidth - 6\tabcolsep) * \real{0.2903}}@{}}
\toprule\noalign{}
\begin{minipage}[b]{\linewidth}\raggedright
Time
\end{minipage} & \begin{minipage}[b]{\linewidth}\raggedright
Frequency
\end{minipage} & \begin{minipage}[b]{\linewidth}\raggedright
Job
\end{minipage} & \begin{minipage}[b]{\linewidth}\raggedright
Details
\end{minipage} \\
\midrule\noalign{}
\endhead
\bottomrule\noalign{}
\endlastfoot
23:00-23:25 & Daily & Agent consolidation & 6 agents, 5-min stagger,
reindex→consolidate \\
23:30 & Daily & Workspace consolidation & Central workspace daily
notes \\
23:35 & Daily & Session wrap-up & SESSION-STATE.md, Notion sync, git
commit \\
04:00 & Weekly (Sun) & Vectorize & Gemini embeddings for new/changed
chunks \\
*/5 min & Continuous & Health check & Watcher heartbeat, service
liveness \\
02:00 & Daily & SQLite backup & Online backup API, 7-day retention
pruning \\
*/6 hours & Continuous & Git backup & Auto-commit memory file changes \\
09:00 & Weekly (Mon) & Token check & Forge CC token verification \\
\end{longtable}
}

\subsubsection{12.2 Backup Strategy}\label{backup-strategy}

Three backup mechanisms operate independently:

\begin{enumerate}
\def\labelenumi{\arabic{enumi}.}
\tightlist
\item
  \textbf{SQLite Online Backup}: Uses better-sqlite3's backup API for
  crash-consistent copies. Daily at 02:00, 7-day retention with
  automatic pruning.
\item
  \textbf{Git Auto-Commit}: Memory directory changes are committed every
  6 hours, providing full change history.
\item
  \textbf{File System}: WAL mode ensures database consistency during
  concurrent reads/writes.
\end{enumerate}

\subsubsection{12.3 LLM Fallback Chain}\label{llm-fallback-chain}

To ensure continuous operation regardless of provider availability:

\textbf{Paid Tier}: Claude Opus → Sonnet → Haiku → GPT-5.1 → Gemini 2.5
\textbf{Free Tier}: Nemotron → Groq Llama70B → Healer → Hunter → Trinity
→ Gemma27B

The fallback is configured in the environment and selected at runtime
based on task complexity and availability.

\begin{center}\rule{0.5\linewidth}{0.5pt}\end{center}

\subsection{13. Dashboard Integration}\label{dashboard-integration}

The TotoClaw Command Center (React 18 + TypeScript + Vite + shadcn/ui)
provides 11 pages including 4 nox-mem-specific views:

\begin{itemize}
\tightlist
\item
  \textbf{Memory Health} (\texttt{/memory}): Real-time system stats,
  vector coverage progress bar, service status indicators, agent
  breakdown table
\item
  \textbf{Knowledge Graph} (\texttt{/knowledge-graph}): Interactive
  force-directed canvas graph, entity type filters, BFS path finder
\item
  \textbf{Agent Intel} (\texttt{/agent-intel}): Agent expertise cards
  with type distribution bars, hybrid search interface, cross-agent
  knowledge entities
\item
  \textbf{System Paper} (\texttt{/system-paper}): Live technical
  analysis with Recharts visualizations (pie, bar, radar, area charts),
  auto-refresh every 60 seconds
\end{itemize}

All data is fetched from the nox-mem API server via TanStack React Query
with configurable polling intervals.

\begin{center}\rule{0.5\linewidth}{0.5pt}\end{center}

\subsection{14. Evolution History}\label{evolution-history}

{\def\LTcaptype{none} % do not increment counter
\begin{longtable}[]{@{}
  >{\raggedright\arraybackslash}p{(\linewidth - 4\tabcolsep) * \real{0.3214}}
  >{\raggedright\arraybackslash}p{(\linewidth - 4\tabcolsep) * \real{0.2143}}
  >{\raggedright\arraybackslash}p{(\linewidth - 4\tabcolsep) * \real{0.4643}}@{}}
\toprule\noalign{}
\begin{minipage}[b]{\linewidth}\raggedright
Version
\end{minipage} & \begin{minipage}[b]{\linewidth}\raggedright
Date
\end{minipage} & \begin{minipage}[b]{\linewidth}\raggedright
Key Changes
\end{minipage} \\
\midrule\noalign{}
\endhead
\bottomrule\noalign{}
\endlastfoot
v1.0 & Mar 14 & SQLite FTS5, basic search, consolidation, Notion sync \\
v2.0 & Mar 17 & MCP server, systemd services, watcher heartbeat,
primer \\
v2.2 & Mar 20 & Cross-agent search, KG v1 (regex), self-improve,
decision versioning \\
v2.5 & Mar 22 & Multi-agent workspace fix (OPENCLAW\_WORKSPACE), gateway
supervision \\
v2.6 & Mar 22 & Hybrid search default (FTS5+Gemini+RRF), 866/866
vectorized \\
v3.0 & Mar 23 & KG v2 (LLM, 384 entities), Cross-Agent Intelligence,
HTTP API, dashboard \\
v3.7 & Apr 23 & Schema V10 (\texttt{retention\_days} v8 + \texttt{pain}
v9 + \texttt{section} v10), entity file format, section\_boost \\
Wave A & May 19 & Additive salience formula, \texttt{tier\_boost}
off-by-default, \texttt{source\_type} backfill (67,949 chunks), G5 V3
ablation (PRs \#150 / \#151 / \#153) \\
G10 Hard Mutex & May 20 & \texttt{section\ ↔\ source\_type} mutex
deployed against \texttt{g9.db} 69,495 chunks (PRs \#181 / \#182) \\
G10d ACTIVE-T2 & May 21 & Conditional mutex gated by
\texttt{query\_entity\_count\ ≤\ 2}, deployed via systemd drop-in
\texttt{NOX\_MUTEX\_QUERY\_ENTITY\_THRESHOLD=2} (PR \#198, decision
D51); multi-hop +1.58\% nDCG / adversarial +3.04\% nDCG recovered \\
F10 Phase A + B & May 21 & Foundation observability dashboards
(\texttt{/observability/health.html} +
\texttt{/observability/evals.html}) deployed via Tailscale tunnel (PRs
\#207 / \#212, decision D53) \\
\end{longtable}
}

\begin{center}\rule{0.5\linewidth}{0.5pt}\end{center}

\subsection{15. Conclusion}\label{conclusion}

nox-mem demonstrates that persistent, searchable, and shareable memory
for AI agent fleets is achievable with commodity infrastructure (single
VPS, SQLite, local LLM). The hybrid search system consistently
outperforms single-method retrieval, particularly for multilingual
content and compound technical terms. The LLM-powered knowledge graph
provides 15x richer entity extraction compared to regex approaches,
while temporal decay ensures the graph stays current without manual
curation. The Wave A empirical evaluation (§5) cravou nDCG@10 = 0.6237
on the entity-flavored golden set (+78.8\% relative over the G3
baseline), with \texttt{section\_boost} identified as the dominant
driver (99.85\% of the lift recovered by A3 alone) and the additive
salience formula validated by the
\texttt{active\ \textgreater{}\ shadow} reversal. The G10d conditional
mutex evolution (§5.5, deployed 2026-05-21) consolida o canonical boost
stack
\texttt{section\_boost\ ×\ source\_type\_boost\ (Hard\ Mutex\ gated\ by\ query\_entity\_count\ ≤\ 2)\ ×\ salience\ v2\ additive}
em produção, recuperando regressões multi-hop e adversarial com diluição
contida em single-hop. A camada F10 (§5.6, decisão D53) torna o estado
de produção verificável a qualquer momento via dashboards Phase A
(\texttt{/observability/health.html}) + Phase B
(\texttt{/observability/evals.html}).

A §6 Q4 COMPARISON está pre-registered
(\texttt{specs/2026-05-23-Q4-comparison-execution-plan.md}) e o
\textbf{smoke de Sat 2026-05-24 15h30 BRT} populou a primeira linha de
§6.3 com números de nox-mem (nDCG@10=0.6380 combined, p50=12ms, gold-hit
13/20 em 20 queries dry-run-sample sobre eval-isolated DB de 5.882
LoCoMo + 940 LongMemEval chunks). O \textbf{run canônico} --- 100
queries × 2 datasets × 6 sistemas (Mem0, Zep, Letta, agentmemory,
EverMind-AI + nox-mem) --- ainda está em execução com 5/6 competitor
adapters em setup, e atualiza as células
\texttt{{[}PENDING\ canonical\ run{]}} quando crava. O gate D43 (top-3
em ≥2 das 4 métricas chave) é avaliado contra o run canônico; o smoke
valida a metodologia + confirma que nox-mem retrieval funciona
end-to-end, destravando a defesa pre-launch da GTM Phase 2.

The cross-agent intelligence layer transforms isolated agent memories
into a collaborative knowledge base, enabling institutional learning
across the fleet. Combined with the live dashboard, the system provides
full observability into the collective memory of the agent organization.

\textbf{Repository:} github.com/totobusnello/nox-workspace
\textbf{Dashboard:} github.com/totobusnello/agent-hub-dashboard
\textbf{Spec:}
Projetos/memoria-nox/specs/2026-03-14-nox-memory-system-design.md

\end{document}
